%%%%%%%%%%%%%%%%%%%%%%%%%%%%%%%%%%%%%%%%%%%%%%%%%%%%%%%%%%%%%%%%%%%%%%
% Lớp 11 | Dãy số. Cấp số cộng. Cấp số nhân | Dãy số
% Dạng 1D2?1-5: Dãy số bị chặn
% Số câu: 76  (sinh tự động bởi tach_cau.py; từ nay thầy sửa trực tiếp file này)
%%%%%%%%%%%%%%%%%%%%%%%%%%%%%%%%%%%%%%%%%%%%%%%%%%%%%%%%%%%%%%%%%%%%%%

%%%%%%%%%%%%%%%% TRẮC NGHIỆM 4 PHƯƠNG ÁN %%%%%%%%%%%%%%%%

% nguồn: Ok_Toan11_TracNghiem_Lan2
\begin{ex}%[1D2H1-5]%[ma:8ac3b4fd]
	Cho dãy số $(u_n)$ với ${u_n}=\dfrac{3n-1}{3n+1}$ . Dãy số $\left(u_n\right)$ bị chặn  trên bởi số nào dưới đây?
	 \choice
		 {$\dfrac{1}{3}$}
		{\True  1}
		{$\dfrac{1}{2}$} 
		 {0}
	 \loigiai{
	Ta có \[u_n=\dfrac{3n-1}{3n+1}=1-\dfrac{2}{3n+1}<1.\] Suy ra dãy $\left(u_n\right)$ bị chặn  trên bởi số $1$.}
\end{ex}

% nguồn: Ok_Toan11_TracNghiem_Lan2
\begin{ex}%[1D2H1-5]%[ma:f5cb6e67]
	Cho dãy số $(u_n)$, với $u_n=(-1)^n$. Mệnh đề  nào sau đây đúng?
	\choice 
		{Dãy  số $(u_n)$ là dãy số tăng}
		{Dãy  số $(u_n)$ là dãy số giảm}
		{\True Dãy số $(u_n)$ là  dãy số bị chặn}
		{Dãy số $(u_n)$  là dãy số không bị chặn}
	 \loigiai{
		Ta có $u_1=-1$; $u_2=1$; $u_3=-1$.\\ 
		 Do đó dãy $(u_n)$ là dãy số không tăng, không giảm.\\
		Lại có $u_n=(-1)^n=1$ khi $n$ chẵn và $u_n=(-1)^n=-1$ với $n$  lẻ.\\
		Do đó $-1\leq u_n \leq 1, \forall n\in  \mathbb{N}^{*}$.\\
		Vậy dãy số $(u_n)$ là dãy số bị chặn. 
	} 
\end{ex}

% nguồn: Ok_Toan11_TracNghiem_Lan2
\begin{ex}%[1D2H1-5]%[ma:de8e4189]
	Trong các dãy  số sau, dãy nào là dãy số bị chặn?
	\choice 
		 {\True $u_n=\dfrac{2n+1}{n+1}$}
		 {$u_n=2n+\sin(n)$}
		{$u_n=n^2$} 
		 {$u_n=n^3-1$}
	 \loigiai{
		Xét dãy số  $u_n=\dfrac{2n+1}{n+1}$, ta có
		\begin{itemize}
			\item $u_n=\dfrac{2n+1}{n+1}>0;\forall  n\in{\mathbb{N}^*}\Rightarrow $ dãy $\left(u_n\right)$ bị chặn dưới bởi giá trị $0$.
			\item $u_n=\dfrac{2n+1}{n+1}=2-\dfrac{1}{n+1}<2;\forall n\in{\mathbb{N}^*}\Rightarrow $ dãy $\left(u_n\right)$ bị  chặn trên bởi $2$.
		\end{itemize}
	Vậy dãy $\left(u_n\right)$  là dãy bị chặn.}
\end{ex}

% nguồn: Ok_Toan11_TracNghiem_Lan2
\begin{ex}%[1D2H1-5]%[ma:ee23b52b]
	Xét tính bị chặn của dãy số  $(u_n)$ với $u_n=4-3n-n^2$.
	 \choice
		{Dãy $(u_n)$ bị chặn} 
		{Dãy $(u_n)$ không bị chặn} 
		{\True Dãy $(u_n)$ chỉ bị chặn  trên}
		 {Dãy $(u_n)$ chỉ bị chặn dưới}
	 \loigiai{
		Vì $u_n=4-3n-n^2=\dfrac{25}{4}-\left(n+\dfrac{3}{2}\right)^2<\dfrac{25}{4}$, $\forall n\in{\mathbb{N}^*}$ nên dãy số $\left(u_n\right)$ bị  chặn trên.\\
	Dãy $\left(u_n\right)$ không bị  chặn dưới.}
\end{ex}

% nguồn: Ok_Toan11_TracNghiem_Lan2
\begin{ex}%[1D2H1-5]%[ma:8c76091a]
	Xét tính tăng, giảm và bị chặn của  dãy số $(u_n)$, biết $u_n=\dfrac{2n-13}{3n-2}$.
	\choice 
		{\True Dãy số tăng, bị  chặn}
		{Dãy  số giảm, bị chặn}
		{Dãy số không tăng, không giảm,  không bị chặn}
		{Dãy số không  tăng, không giảm, bị chặn}
	\loigiai{ 
		 Ta có $u_{n+1}-u_n=\dfrac{2n-11}{3n+1}-\dfrac{2n-13}{3n-2}=\dfrac{34}{(3n+1)(3n-2)} > 0, \forall n\in \mathbb{N}^{*}$.\\
		Suy ra $(u_n)$ là dãy số tăng. Do đó dãy bị chặn  dưới bởi $u_1=-\dfrac{9}{4}$.\\
		Mặt khác $u_n=\dfrac{2}{3}-\dfrac{35}{3(3n-2)} <\dfrac{2}{3}, \forall n\in  \mathbb{N}^{*}$.\\
		Suy ra  $(u_n)$ bị chặn trên.\\
		Vậy dãy số $(u_n)$ là dãy bị chặn. 
	} 
\end{ex}

% nguồn: Ok_Toan11_TracNghiem_Lan2
\begin{ex}%[1D2H1-5]%[ma:47d0244b]
	Trong các dãy số $\left(u_n\right)$ có  số hạng tổng quát $u_n$ dưới đây, dãy số nào là dãy bị chặn?
	\choice 
		{$u_n=\sqrt{n^2+2}$} 
		{\True $u_n=\dfrac{n}{2n+1}$} 
		{$u_n=3^n-1$} 
		 {$u_n=n+\dfrac{2}{n}$}
	 \loigiai{
		Ta có $u_n=\dfrac{n}{2n+1}=\dfrac{1}{2}-\dfrac{1}{2(2n+1)}<\dfrac{1}{2}\Rightarrow $  $u_n<\dfrac{1}{2}$.\\
		Mặt khác ta  có $u_n=\dfrac{n}{2n+1}>0\,\forall n\in\mathbb{N}*$ $\Rightarrow 0<u_n<\dfrac{1}{2}$.\\
	Vậy dãy số $u_n=\dfrac{n}{2n+1}$  bị chặn.}
\end{ex}

% nguồn: Ok_Toan11_TracNghiem_Lan2
\begin{ex}%[1D2H1-5]%[ma:f24acae1]
	Cho dãy $\left(u_n\right)$ với $u_n=\dfrac{n+2018}{2018n+1}$. Chọn  khẳng định đúng trong các khẳng định sau.
	 \choice
		{Dãy  $\left(u_n\right)$ chỉ bị chặn dưới}
		{\True Dãy $\left(u_n\right)$ bị chặn} 
		{Dãy $\left(u_n\right)$ không bị  chặn trên và không bị chặn dưới}
		{Dãy  $\left(u_n\right)$ bị chặn trên nhưng không bị chặn dưới}
	 \loigiai{
		 Ta có \[u_n=\dfrac{n+2018}{2018n+1}=\dfrac{1}{2018}+\dfrac{2017\cdot 2019}{2018\left(2018n+1\right)}.\]
		 Do đó $\left(u_n\right)$ là dãy giảm.\\
		Mặt khác, $u_1=1$ và $\forall  n\in{\mathbb{N}^*},\,u_n>0$ thì $ 0<u_n\le 1.$\\
	Suy ra  dãy $\left(u_n\right)$ bị chặn.}
\end{ex}

% nguồn: Ok_Toan11_TracNghiem_Lan2
\begin{ex}%[1D2H1-5]%[ma:31192ffd]
	Xét tính bị chặn của dãy số $(u_n)$ với $u_n=\dfrac{2n+1}{n+2}$. 
	\choice 
		 {\True Bị chặn}
		{Không bị chặn trên } 
		{Chỉ  bị chặn trên}
		{Chỉ bị  chặn dưới}
	\loigiai{ 
		Ta  có \[u_n=\dfrac{2n+1}{n+2}=2-\dfrac{3}{n+2}\Rightarrow 0<u_n<2,\, \forall n \in \mathbb{N}^*.\]
	 nên dãy $(u_n)$ bị chặn.}
\end{ex}

% nguồn: Ok_Toan11_TracNghiem_Lan2
\begin{ex}%[1D2H1-5]%[ma:2e77c77b]
	Cho dãy số $\left(u_n\right)$ với  $u_n=2+5^{1-n}$ . Kết luận nào sau đây là đúng?
	\choice 
		{Dãy số $\left(u_n\right)$ không đơn điệu} 
		{Dãy số $\left(u_n\right)$ giảm và  không bị chặn}
		{Dãy số  $\left(u_n\right)$ tăng}
		{\True  Dãy số $\left(u_n\right)$ giảm và bị chặn}
	 \loigiai{
		 Ta có
		\begin{eqnarray*}
			 u_{n+1}-u_n&=&\left(2+5^{-n}\right)-\left(2+5^{1-n}\right)\\
			 &=&5^{-n}-5^{1-n}\\
			 &=&\dfrac{1}{5^n}-\dfrac{1}{5^{n-1}}\\
			&=&\dfrac{1}{5^n}-\dfrac{5}{5^n}\\ 
			&=&-\dfrac{4}{5^n}<0,\forall n\in{\mathbb{N}^*}. 
		\end{eqnarray*}
		Vậy dãy số  $\left(u_n\right)$ là dãy số giảm.\\
		Ta  lại có
		$\heva{&u_n=2+5^{1-n}>2,\forall  n\in{\mathbb{N}^*} \\
		 &u_n=2+\dfrac{5}{5^n}\le 3,\forall n\in{\mathbb{N}^*}}$
	nên dãy số $\left(u_n\right)$  là dãy số bị chặn.}
\end{ex}

% nguồn: Ok_Toan11_TracNghiem_Lan2
\begin{ex}%[1D2H1-5]%[ma:9ed937ea]
	Cho dãy số $\left(u_n\right)$, biết $u_n=\cos n+\sin n.$ Dãy số $\left(u_n\right)$ bị chặn  trên bởi số nào dưới đây?
	\choice 
		 {$0$}
		 {$1$}
		 {\True $\sqrt{2}$}
		{$-1$} 
	 \loigiai{
		Ta có \[u_n=\cos n+\sin n=\sqrt{2}\sin\left(n+\dfrac{\pi}{4}\right)\le  \sqrt{2},\, \forall n \in \mathbb{N}^*.\]
		Suy ra  dãy $\left(u_n\right)$ bị chặn trên bởi số $\sqrt{2}$.
	 }
\end{ex}

% nguồn: Ok_Toan11_TracNghiem_Lan2
\begin{ex}%[1D2H1-5]%[ma:69cc2c78]
	Chọn kết luận  {\bf sai}.
	 \choice
		 {\True Dãy số $\left(2n-1\right)$ tăng và bị chặn trên}
		{Dãy số $\left(\dfrac{1}{n+1}\right)$ giảm và  bị chặn dưới}
		{Dãy số $\left(-\dfrac{1}{n}\right)$  tăng và bị chặn trên}
		{Dãy số $\left(\dfrac{1}{3\cdot 2^n}\right)$  giảm và bị chặn dưới}
	 \loigiai{
		\begin{itemize}
			 \item Dãy số $\left(2n-1\right)$ tăng và bị chặn dưới.
			\item  Dãy số $\left(\dfrac{1}{n+1}\right)$ giảm và bị chặn dưới.
			\item Dãy số $\left(-\dfrac{1}{n}\right)$ tăng và bị chặn trên. 
			\item Dãy số $\left(\dfrac{1}{3.2^n}\right)$  giảm và bị chặn dưới
		\end{itemize}
	 }
\end{ex}

% nguồn: Ok_Toan11_TracNghiem_Lan2
\begin{ex}%[1D2H1-5]%[ma:f5921941]
	Xét tính bị chặn của  dãy số sau $u_n=4-3n-n^2$.
	 \choice
		{Bị chặn} 
		{Không  bị chặn}
		{\True  Bị chặn trên}
		{Bị  chặn dưới}
	 \loigiai{
		Ta có $u_{n+1}-u_n=4-3(n+1)-(n+1)^2-(4-3n-n^2)=-2n-4<0, \forall  n\in\mathbb{N}^*$.\\
		Nên  $(u_n)$ là dãy số giảm.\\
		Do  đó $u_n\geq u_1=0$.\\
		Vậy dãy số  $(u_n)$ bị chặn trên.
	 }
\end{ex}

% nguồn: Ok_Toan11_TracNghiem_Lan2
\begin{ex}%[1D2H1-5]%[ma:8efbf0c8]
	 Xét tính bị chặn của các dãy số $u_n=3n-1$.
	 \choice
		{Bị  chặn}
		{Bị chặn trên} 
		{\True  Bị chặn dưới}
		{Không bị  chặn dưới}
	\loigiai{ 
		 Ta có $u_{n+1}-u_n=3(n+1)-1-(3n-1)=3>0, \forall n\in \mathbb{N}^{*}$.\\
		Nên  $(u_n)$ là dãy số tăng.\\
		Do đó $u_n\geq u_1=2, \forall  n\in \mathbb{N}^{*}$.\\
		 Vậy dãy số $(u_n)$ bị chặn dưới.
	 }
\end{ex}

% nguồn: Ok_Toan11_TracNghiem_Lan2
\begin{ex}%[1D2H1-5]%[ma:1d662ccf]
	Xét tính tăng, giảm  và bị chặn của dãy số $(u_n)$, biết $ {u_n}=n^3+2n+1$.
	\choice 
		{Dãy số $(u_n)$  tăng và bị chặn}
		{Dãy số $(u_n)$ giảm  và bị chặn}
		{\True Dãy  số $(u_n)$ tăng và chỉ bị chặn dưới}
		{ Dãy số $(u_n)$ giảm và chỉ bị chặn  trên}
	 \loigiai{
		 Ta có \[u_{n+1}-u_n=(n+1)^3+2(n+1)-n^3-2n=3n^2+3n+3>0,\, \forall n \in \mathbb{N}^*.\]
		Mặt khác $u_n>1,\forall n$ và khi $n$ càng lớn thì $u_n$ càng lớn.\\ 
	Vậy dãy $(u_n)$ là dãy tăng  và chỉ bị chặn dưới.}
\end{ex}

% nguồn: Ok_Toan11_TracNghiem_Lan2
\begin{ex}%[1D2H1-5]%[ma:6245fb3b]
	Trong các dãy số  $(u_n)$ sau, dãy số nào bị chặn?
	\choice 
		 {$u_n=n-\sin 3n$}
		 {$u_n=\dfrac{n^2+1}{n}$}
		{\True $u_n=\dfrac{1}{n\left(n+1\right)}$} 
		{$u_n=n.\sin\left(3n-1\right)$} 
	 \loigiai{
	 Ta có $0<u_n=\dfrac{1}{n\left(n+1\right)}\le\dfrac{1}{2},\,\forall n\in{\mathbb{N}^*}$, suy ra dãy số $(u_n)$ với $u_n=\dfrac{1}{n\left(n+1\right)}$ bị chặn.}
\end{ex}

% nguồn: Ok_Toan11_TracNghiem_Lan2
\begin{ex}%[1D2H1-5]%[ma:f4dcc236]
	Dãy  số nào sau đây là dãy số bị chặn?
	\choice 
		 {$u_n=2n+1$}
		{$u_n=-\dfrac{n^2}{n+1}$} 
		 {$u_n=2^n$}
		 {\True $u_n=\dfrac{n+1}{n}$}
	 \loigiai{
		Xét dãy $(u_n)$  với $u_n=\dfrac{n+1}{n}$.\\
		Ta có \[u_n=\dfrac{n+1}{n}=1+\dfrac{1}{n}>1,\,\forall  n\in{\mathbb{N}^*}.\]
		 Suy ra dãy $(u_n)$ bị chặn dưới.\\
		Lại  có
		\[u_n=1+\dfrac{1}{n}<1+\dfrac{1}{1}=2,\,\forall n\in{\mathbb{N}^*}.\] 
		Suy ra dãy $(u_n)$  bị chặn trên.\\
	Vậy dãy $(u_n)$  bị chặn.}
\end{ex}

% nguồn: Ok_Toan11_TracNghiem_Lan2
\begin{ex}%[1D2H1-5]%[ma:697ba598]
	Cho dãy số $\left(u_n\right)$, biết $u_n=\cos n+\sin n.$ Dãy số $\left(u_n\right)$ bị chặn dưới  bởi số nào dưới đây?
	 \choice
		 {0}
		{$-1$} 
		{\True  $-\sqrt{2}$}
		{$1$} 
	 \loigiai{
		Ta  có \[u_n=\cos n+\sin n=\sqrt{2}\sin\left(n+\dfrac{\pi}{4}\right)\ge -\sqrt{2},\, \forall n \in \mathbb{N}^*.\]
	Suy ra dãy  $\left(u_n\right)$ bị chặn dưới bởi số $-\sqrt{2}$.}
\end{ex}

% nguồn: Ok_Toan11_TracNghiem_Lan2
\begin{ex}%[1D2H1-5]%[ma:68e108c5]
	Xét tính tăng  giảm và bị chặn của dãy số $(u_n)$ với ${u_n}=\dfrac{n+1}{n+2}$.
	\choice 
		{\True Tăng và bị  chặn}
		{Giảm và bị  chặn}
		{Tăng và không  bị chặn}
		{Giảm và không bị chặn} 
	 \loigiai{
		Ta có  \[u_{n+1}-u_n=\dfrac{n+2}{n+3}-\dfrac{n+1}{n+2}=\dfrac{(n+2)^2-(n+3)(n+1)}{(n+2)(n+3)}=\dfrac{1}{(n+2)(n+3)}>0, \forall n \in \mathbb{N}^*.\]
		Và  \[0<u_n=\dfrac{n+1}{n+2}<\dfrac{n+2}{n+2}=1,\,\forall n\in \mathbb{N}^*.\]
	Vậy dãy $(u_n)$ là dãy tăng và  bị chặn.}
\end{ex}

% nguồn: Ok_Toan11_TracNghiem_Lan2
\begin{ex}%[1D2H1-5]%[ma:e07e341e]
	Xét tính bị chặn của các dãy số  sau $u_n=\dfrac{2n+1}{n+2}$.
	 \choice
		 {\True Bị chặn}
		 {Không bị chặn}
		{Bị  chặn trên}
		{Bị  chặn dưới}
	 \loigiai{
		Ta có  $u_{n+1}-u_n=\dfrac{2(n+1)+1}{(n+1)+2}-\dfrac{2n+1}{n+2}=\dfrac{3}{(n+3)(n+2)}>0, \forall n\in \mathbb{N}^{*}$.\\
		 Nên $(u_n)$ là dãy số tăng. Suy ra $u_n\geq u_1=1, \forall n\in \mathbb{N}^{*}$.\\
		Do đó $(u_n)$ bị  chặn dưới.\\
		Lại có $u_n=\dfrac{2n+1}{n+2}=\dfrac{2(n+2)-1}{n+2}=2-\dfrac{1}{n+2}<2, \forall  n\in \mathbb{N}^{*}$.\\
		Do đó $(u_n)$ bị chặn trên.\\ 
		Vậy dãy số $(u_n)$ bị chặn. 
	 }
\end{ex}

% nguồn: Ok_Toan11_TracNghiem_Lan2
\begin{ex}%[1D2H1-5]%[ma:04723729]
	Trong các dãy số $\left(u_n\right)$ cho dưới đây dãy số nào là dãy số bị  chặn?
	 \choice
		 {$u_n=\dfrac{n^3}{n^2+1}$}
		{$u_n=n^2+2017$} 
		 {$u_n=(-1)^n(n+2)$}
		 {\True $u_n=\dfrac{n}{n^2+1}$}
	 \loigiai{
	Ta có $0<u_n=\dfrac{n}{n^2+1}\le\dfrac{1}{2},\,\forall n\in{\mathbb{N}^*}$, suy ra dãy $(u_n)$ với $u_n=\dfrac{n}{n^2+1}$ bị  chặn.}
\end{ex}

% nguồn: Ok_TracNghiemLop11_Lan1
\begin{ex}%[1D2H1-5]%[ma:f9f638d3]
	Xét tính bị chặn của dãy  số $u_n=-n^2+6n-7$.
	\choice 
		 {Bị chặn}
		{Không bị chặn} 
		{\True  Bị chặn trên}
		{Bị chặn dưới} 
	 \loigiai{
		Ta  có $u_n=-n^2+6n-7=2-\left(n+3\right)^2< 7$.
		\begin{itemize}
			\item  Dãy số $\left(u_n\right)$ bị chặn trên;
			\item Dãy dãy  $\left(u_n\right)$ không bị chặn dưới.
		\end{itemize}
	 }
\end{ex}

% nguồn: Ok_TracNghiemLop11_Lan1
\begin{ex}%[1D2H1-5]%[ma:ce7ebe7b]
	Cho dãy số $(u_n)$ với  $u_n=25n^2+10n+9$. Chọn khẳng định đúng
	 \choice
		 {\True Dãy số trên bị chặn dưới}
		{Dãy  số trên bị chặn trên}
		{Dãy  số trên không bị chặn}
		{Dãy số trên bị  chặn}
	\loigiai{ 
		Ta có: $u_n=25 n^2+10 n+9=\left(25 n^2+10 n+1\right)+8=(5 n+1)^2+8 \geq 8$ với  mọi $n \in \mathbb{N}^*$.\\
		Do đó, dãy số $\left(u_n\right)$  bị chặn dưới, không bị chặn trên.
	} 
\end{ex}

% nguồn: Ok_TracNghiemLop11_Lan1
\begin{ex}%[1D2H1-5]%[ma:3b87766a]
	Trong các dãy số $(u_{n})$ sau, dãy số nào là dãy  số bị chặn?
	\choice 
		{\True $u_{n}=\dfrac{2n+1}{n+1}$} 
		{$u_{n} =2n+\sin n$} 
		 {$u_{n} =n^{2} $}
		{$u_{n} =n^{3} -1$} 
	 \loigiai{
		Xét đáp án A, ta có $0<u_{n} =\dfrac{2n+1}{n+1} $$<2, \forall n\in  \mathbb{N}$.
	 }
\end{ex}

% nguồn: Ok_Toan11_TracNghiem_Lan2
\begin{ex}%[1D2V1-5]%[ma:53989bdf]
	Trong các dãy số $(u_n)$ cho bởi số hạng tổng quát $u_n$ sau, dãy số nào bị  chặn?
	 \choice
		{\True  $u_n=\dfrac{1}{2^n}$}
		{$u_n=3^n$} 
		 {$u_n=\sqrt{n+1}$}
		{$u_n=n^2+1$} 
	 \loigiai{
		\begin{itemize}
			\item Xét dãy  $(u_n)$ với $u_n=\dfrac{1}{2^n}$.\\
			Ta có $0< u_n=\dfrac{1}{2^n} \leq \dfrac{1}{2}, \forall n\in  \mathbb{N}^{*}$.\\
			Vậy dãy  số $(u_n)$ bị chặn.
			\item Xét dãy $(u_n)$ với $u_n=3^n$.\\ 
			 Ta có $\dfrac{u_{n+1}}{u_n}=\dfrac{3^{n+1}}{3^n}=3>1, \forall n\in \mathbb{N}^{*}$, nên $(u_n)$ là dãy số tăng.\\
			Do đó  $u_n\geq u_1=3, \forall n\in \mathbb{N}^{*}$.\\
			 Vậy dãy số $(u_n)$ bị chặn dưới.
			\item Xét dãy $(u_n)$ với $u_n=\sqrt{n+1}$.\\ 
			Ta có $\dfrac{u_{n+1}}{u_n}=\sqrt{\dfrac{n+2}{n+1}}=\sqrt{1+\dfrac{1}{n+1}}>1, \forall n\in \mathbb{N}^{*}$, nên $(u_n)$  là dãy số tăng.\\
			Do  đó $u_n\geq u_1=\sqrt2, \forall n\in \mathbb{N}^{*}$.\\
			Vậy dãy số $(u_n)$ bị chặn  dưới.
			\item Xét  dãy $(u_n)$ với $u_n=n^2+1$.\\
			 Ta có $u_{n+1}-u_n=(n+1)^2+1-(n^2+1)=2n+3>0, \forall n\in \mathbb{N}^{*}$, nên $(u_n)$ là dãy số tăng.\\
			Do đó  $u_n\geq u_1=2, \forall n\in \mathbb{N}^{*}$.\\
			Vậy  dãy số $(u_n)$ bị chặn dưới.
		\end{itemize}
	 }
\end{ex}

% nguồn: Ok_Toan11_TracNghiem_Lan2
\begin{ex}%[1D2V1-5]%[ma:7b0f3258]
	Trong các dãy số $(u_n)$ sau,  dãy số nào bị chặn?
	 \choice
		 {$u_n=n+\dfrac{1}{n}$}
		 {$u_n=n+1$}
		{\True  $u_n=\dfrac{n}{2n^2+1}$}
		{$u_n=n^2+n+1$} 
	\loigiai{ 
		\begin{itemize}
			 \item Xét dãy $(u_n)$ với $u_n=n+\dfrac{1}{n}$.\\
			Ta có $1<u_n<n+1, \forall n\in\mathbb{N}^*$,  nên $(u_n)$ bị chặn dưới, nhưng không bị chặn trên.
			\item Xét dãy $(u_n)$  với $u_n=n+1$.\\
			Ta có $2 \leq n+1 \leq 2n, \forall n\in\mathbb{N}^*$, nên $(u_n)$ bị chặn dưới, nhưng không bị chặn trên. 
			\item Xét dãy $(u_n)$ với $u_n=\dfrac{n}{2n^2+1}$.\\ 
			Ta  có $n<2n^2+1, \forall n\in\mathbb{N}^*$ nên $0<\dfrac{n}{2n^2+1}<1, \forall n\in\mathbb{N}^*$.\\
			Vậy dãy  số $(u_n)$ là dãy bị chặn.
			\item Xét dãy $(u_n)$  với $u_n=n^2+n+1$.\\
			Ta có $3\geq n^2+n+1 < 3n^2, \forall n\in\mathbb{N}^*$ nên dãy $(u_n)$ bị chặn dưới,  nhưng không bị chặn trên.
		\end{itemize}
	 }
\end{ex}

% nguồn: Ok_Toan11_TracNghiem_Lan2
\begin{ex}%[1D2V1-5]%[ma:8095e350]
	Xét tính bị chặn của  các dãy số $(u_n)$, biết $\heva{&{u_1}=1\\
	&{u_n}=\dfrac{u_{n-1}+2}{u_{n-1}+1},\,(n\ge 2).} 
	$ 
	\choice 
		{\True Bị  chặn}
		 {Không bị chặn và không bị chặn dưới}
		{Chỉ bị chặn trên} 
		{Chỉ bị chặn  dưới}
	\loigiai{ 
		 Dễ thấy $u_n>0,\,\forall n\in\mathbb{N}^*$.\\
		Ta có \[u_n=\dfrac{u_{n-1}+2}{u_{n-1}+1}  = 1+\dfrac{1}{u_{n-1}+1}\Rightarrow 1<u_n<2, \,\forall n\in\mathbb{N}^*.\]
	 Suy ra dãy $(u_n)$ bị chặn.}
\end{ex}

% nguồn: Ok_Toan11_TracNghiem_Lan2
\begin{ex}%[1D2V1-5]%[ma:32d0a75b]
	Xét dãy số  $(u_n)$ với $u_n=\dfrac{n+1}{\sqrt{n^2+1}}$. Phát biểu nào sau đây đúng nhất?
	\choice 
		{\True Dãy số $(u_n)$ bị chặn} 
		{Dãy số $(u_n)$ không  bị chặn}
		 {Dãy số $(u_n)$ bị chặn trên}
		 {Dãy số $(u_n)$ bị chặn dưới}
	\loigiai{ 
		Ta có $0< u_n=\dfrac{n+1}{\sqrt{n^2+1}}=\sqrt{\dfrac{n^2+2n+1}{n^2+1}}=\sqrt{1+\dfrac{2n}{n^2+1}} \leq \sqrt{1+\dfrac{2n}{2n}}=\sqrt{2} \;, \forall  n\in\mathbb{N}^*$.\\
		 Vậy dãy số $(u_n)$ bị chặn.
	 }
\end{ex}

% nguồn: Ok_Toan11_TracNghiem_Lan2
\begin{ex}%[1D2V1-5]%[ma:def5e999]
	Trong các dãy số $(u_n)$ cho bởi  số hạng tổng quát $(u_n)$ sau, dãy số nào bị chặn?
	 \choice
		{$u_n=n^2$} 
		{$u_n=2^n$} 
		{\True  $u_n=\dfrac{1}{n}$}
		{$u_n=\sqrt{n+1}$} 
	 \loigiai{
		\begin{itemize}
			 \item Xét dãy $(u_n)$ với $u_n=n^2$.\\
			Ta có $u_{n+1}-u_n=2n+1>0, \forall n\in \mathbb{N}^{*}$, nên $(u_n)$ là dãy số  tăng.\\
			Do đó $u_n\geq  u_1=1, \forall n\in \mathbb{N}^{*}$.\\
			Vậy dãy số $(u_n)$ bị chặn  dưới.
			\item Xét dãy $(u_n)$ với  $u_n=2^n$.\\
			Ta có $\dfrac{u_{n+1}}{u_n}=\dfrac{2^{n+1}}{2^n}=2>1, \forall n\in \mathbb{N}^{*}$, nên $(u_n)$ là dãy số tăng.\\ 
			Do đó $u_n\geq u_1=2, \forall n\in \mathbb{N}^{*}$.\\ 
			Vậy dãy số $(u_n)$ bị chặn dưới. 
			\item Xét dãy $(u_n)$  với $u_n=\dfrac{1}{n}$.\\
			Ta có  $0< u_n=\dfrac{1}{n} \leq 1, \forall n\in \mathbb{N}^{*}$.\\
			 Vậy dãy số $(u_n)$ bị chặn.
			 \item Xét dãy $(u_n)$ với $u_n=\sqrt{n+1}$.\\
			Ta có $\dfrac{u_{n+1}}{u_n}=\sqrt{\dfrac{n+2}{n+1}}=\sqrt{1+\dfrac{1}{n+1}}>1, \forall n\in \mathbb{N}^{*}$, nên $(u_n)$ là dãy  số tăng.\\
			 Do đó $u_n\geq u_1=\sqrt2, \forall n\in \mathbb{N}^{*}$.\\
			Vậy dãy số $(u_n)$ bị  chặn dưới.
		\end{itemize}
	 }
\end{ex}

% nguồn: Ok_Toan11_TracNghiem_Lan2
\begin{ex}%[1D2V1-5]%[ma:9dae2a83]
	Cho dãy số $(u_n)$ gồm  các số hạng là các số nguyên, thỏa mãn $\log^2u_1+\log{u_1}-6=0$ và $u_{n+1}=u_n+5$, với mọi $n\ge 1,\,n\in \mathbb{N}$. Giá trị lớn nhất của $n$ để $u_n<500$ bằng
	\choice 
		{\True $80$} 
		{$100$} 
		 {$99$}
		{$82$} 
	\loigiai{ 
		\begin{itemize}
			\item  $\log^2u_1+\log{u_1}-6=0\Leftrightarrow \hoac{&\log{u_1}=-3\\
			 &
			\log{u_1}=2}\Leftrightarrow  \hoac{&{u_1}=0,001\, \text{(loại)}\\
			& 
			{u_1}=100.} 
			 $
			\item Từ giả thiết suy ra $(u_n)$  là cấp số cộng có công sai $d=5$.\\
			Do đó, ta  có $u_n=u_1+(n-1)d$.
		\end{itemize}
		Vậy \[u_n<500\Leftrightarrow 5n+95<500 \Leftrightarrow 
		n<81.\] 
	Vậy số $n$  lớn nhất để $u_n<500$ là $80$.}
\end{ex}

% nguồn: Ok_Toan11_TracNghiem_Lan2
\begin{ex}%[1D2V1-5]%[ma:39461fad]
	Xét tính bị chặn của các dãy số $(u_n)$, biết  $u_n=\dfrac{1}{1\cdot 3}+\dfrac{1}{3\cdot 5}+\ldots +\dfrac{1}{\left(2n-1\right)\left(2n+1\right)}$
	 \choice
		{\True Bị chặn} 
		{Không bị chặn trên và không bị  chặn dưới}
		{Chỉ bị  chặn trên}
		{Chỉ bị chặn  dưới}
	\loigiai{ 
		 Rõ ràng $u_n>0,\forall n\in\mathbb{N}*$ nên $\left(u_n\right)$ bị chặn dưới.\\
		Lại  có \[\dfrac{1}{\left(2k-1\right)\left(2k+1\right)}=\dfrac{1}{2}\left(\dfrac{1}{2k-1}-\dfrac{1}{2k+1}\right).\]
		Suy  ra
		\begin{eqnarray*}
			u_n&=&\dfrac{1}{2}\left[\left(1-\dfrac{1}{3}\right)+\left(\dfrac{1}{3}-\dfrac{1}{5}\right)+  \ldots +\left(\dfrac{1}{2n-1}-\dfrac{1}{2n+1}\right)\right]\\
			&=&\dfrac{1}{2}\left(1-\dfrac{1}{2n+1}\right)\\ 
			&<&\dfrac{1}{2},\,\forall n\in\mathbb{N}^* 
		\end{eqnarray*}
		nên $\left(u_n\right)$ bị chặn trên.\\ 
	Kết  luận $\left(u_n\right)$ bị chặn.}
\end{ex}

% nguồn: Ok_Toan11_TracNghiem_Lan2
\begin{ex}%[1D2V1-5]%[ma:6072d63c]
	Xét tính bị  chặn của dãy số $(u_n)$, biết $u_n=\dfrac{1}{1\cdot 3}+\dfrac{1}{2\cdot 4}+\ldots+\dfrac{1}{n\cdot (n+2)}$.
	\choice 
		{\True Bị  chặn}
		{Không bị chặn trên và không bị  chặn dưới}
		{Chỉ bị chặn trên} 
		{Chỉ bị chặn dưới} 
	\loigiai{ 
		Ta có \[0<u_n<\dfrac{1}{1\cdot  2}+\dfrac{1}{2\cdot 3}+\ldots+\dfrac{1}{n\cdot (n+1)}=1-\dfrac{1}{n+1}<1,\,\forall n\in\mathbb{N}^*.\]
	Dãy $(u_n)$ bị  chặn.}
\end{ex}

% nguồn: Ok_Toan11_TracNghiem_Lan2
\begin{ex}%[1D2V1-5]%[ma:b86180fc]
	Trong các dãy số sau đây, dãy số  nào bị chặn?
	 \choice
		 {$u_n=2n$}
		{\True  $u_n=\dfrac{2n+1}{n+2}$}
		 {$u_n=\dfrac{2n+1}{3}$}
		{$u_n=3-2n$} 
	 \loigiai{
		\begin{itemize}
			\item  Xét dãy $(u_n)$ với $u_n=2n$.\\
			Ta có  $u_{n+1}-u_n=2(n+1)-2n=2>0, \forall n\in \mathbb{N}^{*}$, nên $(u_n)$ là dãy số tăng.\\
			Do đó $u_n\geq u_1=2,  \forall n\in \mathbb{N}^{*}$.\\
			Vậy dãy  số $(u_n)$ bị chặn dưới.
			\item Xét  dãy $(u_n)$ với $u_n=\dfrac{2n+1}{n+2}$.\\
			Ta có $u_{n+1}-u_n=\dfrac{2(n+1)+1}{(2n+1)+2}-\dfrac{2n+1}{n+2}=\dfrac{3}{(n+3)(n+2)}>0, \forall n\in  \mathbb{N}^{*}$.\\
			Nên  dãy $(u_n)$ là dãy số tăng.\\
			Do đó $u_n\geq  u_1=1$, suy ra $(u_n)$ bị chặn dưới.\\
			Lại có $u_n=\dfrac{2n+1}{n+2}=2-\dfrac{3}{n+2}<2, \forall  n\in \mathbb{N}^{*}$, suy ra $(u_n)$ bị chặn trên.\\
			Vậy dãy  số $(u_n)$ bị chặn.
			\item Xét dãy $(u_n)$  với $u_n=\dfrac{2n+1}{3}$.\\
			 Ta có $u_{n+1}-u_n=\dfrac{2(n+1)+1}{3}-\dfrac{2n+1}{3}=\dfrac{2}{3}>0, \forall n\in \mathbb{N}^{*}$, nên $(u_n)$ là dãy số tăng.\\
			Do đó $u_n\geq u_1=1, \forall  n\in \mathbb{N}^{*}$.\\
			Vậy dãy số $(u_n)$  bị chặn dưới.
			\item Xét dãy $(u_n)$  với $u_n=3-2n$.\\
			Ta có $u_{n+1}-u_n=-2<0, \forall n\in  \mathbb{N}^{*}$, nên $(u_n)$ là dãy số giảm.\\
			Do  đó $u_n\leq u_1=1, \forall n\in \mathbb{N}^{*}$.\\
			Vậy dãy số $(u_n)$ bị  chặn trên.
		\end{itemize}
	} 
\end{ex}

% nguồn: Ok_Toan11_TracNghiem_Lan2
\begin{ex}%[1D2V1-5]%[ma:bd07b9c0]
	Cho dãy  số $(u_n)$ biết $u_n=\dfrac{1}{2}+\dfrac{1}{2^2}+\dfrac{1}{3^2}+...+\dfrac{1}{n^2}$ . Mệnh đề nào sau đây đúng?
	\choice 
		{Dãy số chỉ bị chặn  dưới}
		{Dãy số  chỉ bị chặn trên}
		{\True Dãy số bị chặn} 
		{Dãy số không bị chặn trên và không  bị chặn dưới}
	 \loigiai{
		Ta  có \[\dfrac{1}{k^2}<\dfrac{1}{\left(k-1\right)k}=\dfrac{1}{k-1}-\dfrac{1}{k},\forall k\ge 2.\]
		Suy ra 
		\begin{eqnarray*}
			 u_n&<&\dfrac{1}{2}+\left(1-\dfrac{1}{2}\right)+\left(\dfrac{1}{2}-\dfrac{1}{3}\right)+\left(\dfrac{1}{3}-\dfrac{1}{4}\right)+\left(\dfrac{1}{5}-\dfrac{1}{6}\right)+\ldots+\left(\dfrac{1}{n-1}-\dfrac{1}{n}\right)\\
			&=&\dfrac{3}{2}-\dfrac{1}{n}\\ 
			&<&\dfrac{3}{2},\, \forall  n\in{\mathbb{N}^*}.
		\end{eqnarray*}
		$\Rightarrow 0<u_n<\dfrac{3}{2},\forall  n\in\mathbb{N}*$.\\
	Vậy  $(u_n)$ bị chặn.}
\end{ex}

% nguồn: Ok_Toan11_TracNghiem_Lan2
\begin{ex}%[1D2V1-5]%[ma:85f0401e]
	Xét tính tăng giảm và  bị chặn của dãy số $(u_n)$, biết $\heva{&{u_1}=2\\
	&{u_{n+1}}=\dfrac{u_n+1}{2},\forall  n\ge 2.}$
	\choice 
		{Dãy số $(u_n)$ tăng và bị chặn} 
		{\True Dãy số  $(u_n)$ giảm và bị chặn}
		{Dãy  số $(u_n)$ tăng và chỉ bị chặn dưới}
		{Dãy số $(u_n)$ giảm và chỉ chặn trên} 
	 \loigiai{
		Dự đoán $1<u_n\le 2,\forall n \in\mathbb{N}^*$.\\ 
		Ta chứng minh  bằng quy nạp.\\
		 Điều này đúng với $n=1$, giả sử $1<u_k\le 2$ với $k\ge 1$, ta có\\
		 $1<u_{k+1}=\dfrac{u_k+1}{2}\le 2$ nên ta có điều phải chứng minh.\\
		Mà $u_{n+1}-u_n=\dfrac{1-u_n}{2}<0,\forall n\in\mathbb{N}^*$.\\ 
	Vậy dãy $(u_n)$ là dãy giảm và bị  chặn.}
\end{ex}

% nguồn: Ok_Toan11_TracNghiem_Lan2
\begin{ex}%[1D2V1-5]%[ma:233f7845]
	Cho  dãy số $(u_n)$ xác định bởi $\heva{&{u_1}=1\\
	 &{u_{n+1}}=u_n+n^3,\forall n\in{\mathbb{N}^*}}$. Tìm số nguyên dương $n$ nhỏ nhất sao cho $\sqrt{u_n-1}\ge 2039190$.
	\choice 
		{$n=2017$} 
		 {$n=2019$}
		{\True $n=2020$} 
		{$n=2018$} 
	\loigiai{ 
		Ta  có
		\[\heva{&u_{n}-u_{n-1}=(n-1)^3  \\
		&u_{n-1}-u_{n-2}=(n-2)^3\\ 
		&\ldots\\ 
		& u_2-u_1=1^3}\Rightarrow u_n-u_1  =1^3+2^3+\ldots (n-1)^3.\]
		Suy  ra $u_n=1+\dfrac{n^2(n-1)^2}{4}\Rightarrow\sqrt{u_n-1}=\dfrac{n(n-1)}{2}$.\\
	Vì vậy  $\sqrt{u_n-1}\ge 2039190 \Rightarrow n\ge 2020$.}
\end{ex}

% nguồn: Ok_TracNghiemLop11_Lan1
\begin{ex}%[1D2V1-5]%[ma:6dd3dd46]
	Cho dãy số $(u_{n})$ với $u_{n} =\dfrac{1}{n+2}$. Trong các mệnh đề dưới  đây, mệnh đề nào đúng?
	 \choice
		{\True Dãy số $(u_{n})$ là dãy số  giảm và bị chặn}
		{Dãy số $(u_{n})$ là dãy số  tăng và bị chặn trên}
		{Dãy  số $(u_{n})$ là dãy số giảm và không bị chặn dưới}
		{Dãy số $(u_{n})$  là dãy số tăng và không bị chặn trên}
	\loigiai{ 
		Ta có $u_{n} =\dfrac{1}{n+2} \Rightarrow u_{n+1} =\dfrac{1}{\left(n+1\right)+2}  =\dfrac{1}{n+3}$.\\
		Suy  ra $u_{n+1} <u_{n}, \forall n \in \mathbb{N}^{*} $.\\
		 Vậy dãy số $(u_{n})$ là dãy số giảm (loại B, D).\\
		Với $\forall n \in {\rm N}^{*}\mathbb{N}^{*} $ ta có $u_{n} =\dfrac{1}{n+2} >0$ và $n+2>1 \Rightarrow u_{n} =\dfrac{1}{n+2} <1$.\\ 
		Vậy dãy số $(u_{n})$ là  dãy số bị chặn.
	} 
\end{ex}

% nguồn: Ok_TracNghiemLop11_Lan1
\begin{ex}%[1D2V1-5]%[ma:f296c682]
	Xét các  câu sau:
	\begin{itemize}
		\item[(1)] Dãy  $1,2,3,\ldots,n,\ldots$ là dãy bị chặn.
		 \item[(2)] Dãy $1,\dfrac{1}{3}$ , $\dfrac{1}{5}$, $\dfrac{1}{7}$, $\ldots$, $\dfrac{1}{2n-1}$, $\ldots$ là dãy bị chặn trên nhưng không bị chặn dưới.
	\end{itemize}
	Mệnh đề nào sau  đây là đúng?
	 \choice
		{Chỉ  có $(2)$ đúng}
		{Chỉ  có $(1)$ đúng}
		 {Cả hai câu đều đúng}
		{\True Cả hai  câu đều sai}
	\loigiai{ 
		Dãy $1,2,3,\ldots,n,\ldots$ là  dãy bị chặn dưới, không bị chặn trên nên không phải dãy số bị chặn.\\
		Dãy $1,\dfrac{1}{3}$  , $\dfrac{1}{5}$, $\dfrac{1}{7}$, $\ldots$, $\dfrac{1}{2n-1}$, $\ldots$ là dãy bị chặn trên tại$1$ và bị chặn dưới tại $0$.\\
	Do đó cả hai câu trên đều  sai.}
\end{ex}

% nguồn: Ok_TracNghiemLop11_Lan1
\begin{ex}%[1D2V1-5]%[ma:e2a14dc5]
	Cho dãy số $\left(u_{n} \right)$ với $u_{n}  =2n-1$. Dãy số $\left(u_{n} \right)$ là dãy số
	 \choice
		{Bị chặn  trên bởi 1}
		 {Giảm}
		{Bị chặn dưới bởi 2} 
		{\True  Tăng}
	 \loigiai{
		$\forall n\in \mathbb{N}^{*}$, ta có $u_{n+1} -u_{n}  =2\left(n+1\right)-1-\left(2n-1\right)=2>0$ nên $u_{n+1} >u_{n}$.\\
	Vậy dãy số $\left(u_{n} \right)$ là dãy tăng.} 
\end{ex}

% nguồn: Ok_TracNghiemLop11_Lan1
\begin{ex}%[1D2V1-5]%[ma:708b6a35]
	Cho dãy số $\left(u_{n} \right)$ có $u_{n} =\dfrac{1}{n\left(n+1\right)}$.  Tìm mệnh đề đúng.
	 \choice
		{Dãy  số $\left(u_{n} \right)$ chỉ bị chặn dưới}
		{Dãy số $\left(u_{n} \right)$ tăng} 
		 {\True Dãy số $\left(u_{n} \right)$ bị chặn}
		{Dãy số $\left(u_{n} \right)$ chỉ bị chặn  trên}
	 \loigiai{
		Xét $\dfrac{u_{n+1} }{u_{n} } =\dfrac{n}{n+2}$.\\ 
		Với $\forall  n\geq$ ta có $0<u_{n} <1$.\\
	Vậy dãy $\left(u_{n} \right)$ bị  chặn.}
\end{ex}

% nguồn: Ok_TracNghiemLop11_Lan1
\begin{ex}%[1D2V1-5]%[ma:4c11257e]
	Cho dãy số $\left(u_{n} \right)$ với $u_{n} =\left(-1\right)^{n}  \sqrt{n} $. Mệnh đề nào sau đây đúng?
	\choice 
		{Dãy số $\left(u_{n}  \right)$ là dãy số bị chặn}
		{Dãy số $\left(u_{n} \right)$ là dãy số  giảm}
		{Dãy số $\left(u_{n}  \right)$ là dãy số tăng}
		{\True Dãy số $\left(u_{n} \right)$ là dãy số không bị chặn} 
	 \loigiai{
	Dãy  số $\left(u_{n} \right)$ là dãy số không bị chặn vì $\lim \left|u_{n} \right|=\lim \sqrt{n} =+\infty $.}
\end{ex}

% nguồn: Ok_TracNghiemLop11_Lan1
\begin{ex}%[1D2V1-5]%[ma:8744422e]
	Cho dãy số $\left(u_{n} \right)$ với  $u_{n} =\dfrac{1}{n^{2} +n} $. Khẳng định nào sau đây là sai?
	\choice 
		{Năm số hạng đầu của  dãy là $\dfrac{1}{2}$, $\dfrac{1}{6};\dfrac{1}{12}$, $\dfrac{1}{20};\dfrac{1}{30}$}
		{\True $\left(u_{n} \right)$ là  dãy số tăng}
		{Bị chặn trên bởi  số $M=\dfrac{1}{2}$}
		{Không bị chặn} 
	\loigiai{ 
		 Ta có
		\begin{eqnarray*}
			& u_{n+1} -u_{n} & =\dfrac{1}{\left(n+1\right)^{2} +\left(n+1\right)}  - \dfrac{1}{n^{2} +n}\\
			&  &= \dfrac{1}{\left(n+1\right)\left(n+2\right)} -\dfrac{1}{n\left(n+1\right)}\\
			&&= \dfrac{-2}{n\left(n+1\right)\left(n+2\right)}  <0.
		\end{eqnarray*}
		với $n\geq  1$.\\
	Do  đó là dãy giảm.}
\end{ex}


%%%%%%%%%%%%%%%% ĐÚNG/SAI %%%%%%%%%%%%%%%%

% nguồn: Ok_dungsai_lop11_Lan1
\begin{ex}%[1D2H1-5]%[ma:243f5771]
 Cho dãy số $(u_n)$ biết $u_n=2^n$. Các mệnh đề sau  đúng hay sai?
 \choiceTF 
  {\True Dãy số  tăng}
  {\True Dãy số không bị chặn trên} 
  {\True Dạy  số bị chặn dưới}
  {\True Số hạng tổng quát của dãy số  là $u_{n+2}=4\cdot 2^{n}$}
 \loigiai{ 
  \begin{itemchoice}
   \itemch Mệnh đề \lq\lq Dãy số tăng\rq\rq\, đúng. 
   \itemch Mệnh đề \lq\lq Dãy số không  bị chặn trên\rq\rq\, đúng.
   \itemch  Mệnh đề \lq\lq Dạy số bị chặn dưới\rq\rq\, đúng.
   \itemch Mệnh đề \lq\lq  Số hạng tổng quát của dãy số là $u_{n+2}=4\cdot 2^{n}$\rq\rq\, đúng.
  \end{itemchoice}
  }
\end{ex}

% nguồn: Ok_dungsai_lop11_Lan1
\begin{ex}%[1D2H1-5]%[ma:e42667f5]
 Cho dãy số $\left( u_n\right)$ với $u_n=\dfrac{an-1}{n+2} $. Các mệnh đề sau đúng  hay sai?
 \choiceTF 
  {\True Dãy số tăng khi $ a>-\dfrac{1}{2} $} 
  {\True Dãy  số giảm khi $ a=0 $}
  {\True  Dãy số bị chặn}
  {\True Dãy số không tăng không giảm khi  $ a=-\dfrac{1}{2} $}
  \loigiai{
  Ta có $ u_{n+1}-u_n=\dfrac{an+a-1}{n+3}-\dfrac{an-1}{n+2}=\dfrac{2a+1}{(n+2)(n+3)} $.\\ 
  Để dãy số  $ (u_n) $ tăng $ \Leftrightarrow u_{n+1}-u_n>0,\forall n\in\mathbb{N}^*\Leftrightarrow 2a+1>0\Leftrightarrow a>-\dfrac{1}{2} $.
  \begin{itemchoice}
   \itemch Mệnh đề \lq\lq Dãy số tăng khi $ a>-\dfrac{1}{2} $\rq\rq\, đúng. 
   \itemch Mệnh đề \lq\lq Dãy số giảm khi $ a=0  $\rq\rq\, đúng.
   \itemch Mệnh đề \lq\lq Dãy số  bị chặn\rq\rq\, đúng.
   \itemch Mệnh đề \lq\lq Dãy số không tăng không  giảm khi $ a=-\dfrac{1}{2} $\rq\rq\, đúng.
  \end{itemchoice}
 } 
\end{ex}

% nguồn: Ok_dungsai_lop11_Lan1
\begin{ex}%[1D2H1-5]%[ma:fef2a5b6]
 Cho dãy số $\left( u_n\right) $ có số  hạng đầu $u_1=1$, $u_n=u_{n-1}+2$, $n\in \mathbb{N},\ n>1$. Các mệnh đề sau đúng hay sai?
  \choiceTF
   {\True $u_2=3$}
   {$u_5=11$}
  {\True Dãy số  $\left( u_n\right) $ là dãy số tăng}
  {Dãy số $\left(  u_n\right) $ là một dãy bị chặn}
 \loigiai{ 
  \begin{itemchoice}
   \itemch Ta có $u_2 =  u_1 + 2 = 1 +2 = 3$.
   \itemch $u_2 = 3$ suy  ra $u_3 = u_2 +2 = 5$ suy ra $u_4 = u_3 + 2 = 7$ suy ra $u_5 = u_4 + 2 = 9$.
   \itemch Vì $\forall n \ge 1$ ta có $u_n -  u_{n-1} =2 >0$ nên dãy số $\left( u_n\right) $ là dãy số tăng.
   \itemch Vì dãy số $\left( u_n\right) $ không bị chặn trên nên dãy  số đã cho không bị chặn.
  \end{itemchoice}
  }
\end{ex}

% nguồn: Ok_dungsai_lop11_Lan1
\begin{ex}%[1D2H1-5]%[ma:f06a5eed]
 Cho dãy số $\left( u_{n} \right)$ có số hạng tổng quát $u_{n}=\dfrac{2n+1}{n+2}$. Khi đó 
  \choiceTF
  {\True Số hạng thứ 100 của dãy số  trên là $u_{100}=\dfrac{67}{34}$}
  {Số $\dfrac{67}{34}$  là số hạng thứ 251 của dãy số}
  {Dãy số  $\left( u_{n} \right)$ là dãy số giảm }
  {\True Dãy số $\left( u_{n} \right)$ là dãy số  bị chặn}
 \loigiai{ 
  \begin{itemchoice}
    \itemch \textbf{Đúng}.\\
   Số  hạng thứ 100 là $u_{100}=\dfrac{2\cdot100+1}{100+2}=\dfrac{67}{34}$.
    \itemch \textbf{Sai}.\\
   Giả sử $u_{n}=\dfrac{167}{84}\Rightarrow \dfrac{2n+1}{n+2}=\dfrac{167}{84}\Rightarrow  n=250$. \\
   Vậy $\dfrac{167}{84}$  là số hạng thứ 250 của dãy số.
   \itemch \textbf{Sai}.\\ 
    Ta có $u_{n+1}=\dfrac{2\left( n+1 \right)+1}{n+1+2}=\dfrac{2n+3}{n+3}$. \\
   $u_{n+1}-u_{n}=\dfrac{2n+3}{n+3}-\dfrac{2n+1}{n+2}=\dfrac{3}{\left( n+3  \right)\left( n+2 \right)}>0,\forall n\in \mathbb{N^{*}}$. \\
   Vậy  dãy số $\left( u_{n} \right)$ là dãy số tăng.
   \itemch  \textbf{Đúng}.\\
   Ta có $u_{n}>0,\forall n\geqslant  1$. Suy ra, dãy số bị chặn dưới. \\
   Mặt khác $u_{n}=\dfrac{2n+1}{n+2}=2-\dfrac{3}{n+2}<2,\forall  n\geqslant 1$. Suy ra dãy số bị chặn trên. \\
   Vậy dãy  số trên bị chặn.
  \end{itemchoice}
 } 
\end{ex}

% nguồn: Ok_dungsai_lop11_Lan1
\begin{ex}%[1D2H1-5]%[ma:f5caf8b9]
 Cho dãy số $(u_n)$ được xác định bởi $\heva{&u_1=3\\&u_{n+1}=2u_n \quad (n \ge 1)}$. Khi  đó
 \choiceTF 
  {Dãy số $(u_n)$ là dãy  số giảm}
  {Dãy số $(u_n)$ là dãy số bị chặn} 
  {\True $u_2  = 6$}
  {\True Công thức tổng quát của $(u_n)$  là $u_n=2^{n-1}\cdot 3$}
  \loigiai{
  \begin{itemchoice}
   \itemch Sai. Dãy số $(u_n)$ là dãy số tăng. Vì $u_n  < u_{n+1},\ \forall n \in \mathbb{N}^*$.
   \itemch Sai. Vì khi $n$ tiến ra $+\infty$ thì $u_n = \left(2^{n-1}\cdot 3\right)$ cũng tiến ra $+\infty\  \forall n \in \mathbb{N}^{\ast}$. Do đó dãy số $(u_n)$ không có chặn trên. Nên $(u_n)$ không bị chặn.
   \itemch Đúng. Vì  $u_2 = 2u_1=2\cdot 3=6$.
   \itemch  Đúng. \\
   \textbf{Cách 1:  Dùng quy nạp.}\\
   Ta  có $u_1 = 3$.\\
   Giả sử $u_k = 2^{k-1}\cdot 3$, đúng $\forall k \ge  1$.\\
   Ta  chứng minh $u_{k+1} = 2^{k}\cdot 3$.
   Thật vậy $u_{k+1} = 2u_k \Rightarrow 2u_k = 2^k \cdot 3 \Rightarrow u_k  = 2^{k-1}\cdot 3$.\\
   Vậy $u_n = 2^{n-1}\cdot  3$, $\forall n \ge 1$.\\
    \textbf{Cách 2: Không dùng quy nạp.}\\
   Dự đoán công thức  tổng quát của dãy $(u_n)$ là $u_n = 2^{n-1}\cdot 3$.\\
   Ta  kiểm tra lại:
   \begin{itemize}
    \item $u_1 = 2^{1-1}\cdot 3  = 2^0\cdot 3 = 3$.
    \item Thay $n = k$ ta được $u_k = 2^{k-1}\cdot 3$. Thay $k = n+1$ ta được $u_{n+1}  = u^n \cdot 3 = 2u^{n-1}\cdot 3 = 2u_n$.
   \end{itemize}
  \end{itemchoice}
 } 
\end{ex}

% nguồn: Ok_dungsai_lop11_Lan1
\begin{ex}%[1D2H1-5]%[ma:99a93fe3]
 Các mệnh đề sau đúng hay sai? 
  \choiceTF
  {\True Dãy số  $u_n=n$ tăng}
  {\True  Dãy số $u_n=(-1)^nn$ không bị chặn}
  {\True Dãy số $u_n=n$ bị chặn dưới không bị  chặn trên}
  {\True  Dãy số $u_n=(-1)^n \dfrac{1}{n}$ dãy số bị chặn}
 \loigiai{ 
  \begin{itemchoice}
   \itemch Mệnh đề \lq\lq Dãy số $u_n=n$ tăng\rq\rq\,  đúng.
   \itemch Mệnh đề \lq\lq Dãy số $u_n=(-1)^nn$  không bị chặn\rq\rq\, đúng.
   \itemch Mệnh đề  \lq\lq Dãy số $u_n=n$ bị chặn dưới không bị chặn trên\rq\rq\, đúng.
   \itemch Mệnh đề \lq\lq Dãy số $u_n=(-1)^n \dfrac{1}{n}$  dãy số bị chặn\rq\rq\, đúng.
  \end{itemchoice}
  }
\end{ex}

% nguồn: Ok_dungsai_lop11_Lan1
\begin{ex}%[1D2H1-5]%[ma:ac5ee455]
 Cho dãy số $(u_n)$ với $u_n=\dfrac{(-1)^{n-1}}{n+1}$.  Các mệnh đề sau đúng hay sai?
 \choiceTF 
  {\True  Số hạng thứ $10$ của dãy số là $\dfrac{-1}{11}$}
  {\True Số hạng thứ $9$  của dãy số là $\dfrac{1}{10}$}
  {Dãy số $(u_n)$ là  một dãy số giảm}
  {\True  Dãy số $(u_n)$ bị chặn}
 \loigiai{ 
  Dễ  thấy $\left| u_n\right|=\left| \dfrac{(-1)^{n-1}}{n+1}\right|=\dfrac{1}{n+1}<1$, $\forall n\in \mathbb{N}^*$ nên $(u_n)$ là dãy số bị chặn.\\
  Lại có $u_9=\dfrac{1}{10}$; $u_{10}=\dfrac{-1}{11}$; $u_{11}=\dfrac{1}{12}$;  $u_{12}=\dfrac{-1}{13}$ suy ra dãy $(u_n)$ không phải là dãy số tăng cũng không phải là dãy số giảm. \\
  Do đó đáp án \lq\lq Dãy số $(u_n)$ là một  dãy số giảm\rq\rq\, sai.
  \begin{itemchoice}
   \itemch Mệnh đề \lq\lq Số hạng thứ $10$  của dãy số là $\dfrac{-1}{11}$\rq\rq\, đúng.
   \itemch Mệnh đề \lq\lq Số hạng thứ $9$  của dãy số là $\dfrac{1}{10}$\rq\rq\, đúng.
   \itemch Mệnh đề \lq\lq Dãy số $(u_n)$ là một dãy số  giảm\rq\rq\, sai.
   \itemch Mệnh đề \lq\lq Dãy  số $(u_n)$ bị chặn\rq\rq\, đúng.
  \end{itemchoice}
  }
\end{ex}

% nguồn: Ok_dungsai_lop11_Lan1
\begin{ex}%[1D2H1-5]%[ma:165040d3]
 Các mệnh  đề sau đúng hay sai?
 \choiceTF 
  {\True Cho dãy số $(u_n)$, với $u_n=5^{n+1}$. Khi  đó số hạng $u_{n-1}=5^n$}
  {Dãy số $u_n=\dfrac{\sqrt{2}}{2^n}$ là một dãy  số tăng}
  {Dãy  số $u_n=-\sin n$ là một dãy số tăng}
  {\True Dãy số  $u_n=\dfrac{3\sin n+4\cos n}{n+1}$ là một dãy số bị chặn}
  \loigiai{
  \begin{itemchoice}
   \itemch Mệnh đề \lq\lq Cho dãy số  $(u_n),$ với $u_n=5^{n+1}$. Khi đó số hạng $u_{n-1}=5^n$\rq\rq\, đúng.
   \itemch Mệnh đề \lq\lq  Dãy số $u_n=\dfrac{\sqrt{2}}{2^n}$ là một dãy số tăng\rq\rq\, sai.
   \itemch  Mệnh đề \lq\lq Dãy số $u_n=-\sin n$ là một dãy số tăng\rq\rq\, sai.
   \itemch Mệnh đề \lq\lq Dãy số $u_n=\dfrac{3\sin n+4\cos n}{n+1}$ là một dãy  số bị chặn\rq\rq\, đúng.
  \end{itemchoice}
  }
\end{ex}

% nguồn: Ok_dungsai_lop11_Lan1
\begin{ex}%[1D2H1-5]%[ma:6fb359a3]
  Các mệnh đề sau đúng hay sai?
  \choiceTF
  {\True  Cho dãy số $(u_n)$ với $u_n = (-1)^n \cdot \dfrac{n}{n+1}$. Khi đó $u_8=\dfrac{8}{9}$}
  {\True Dãy số $\left(u_n\right)$ biết $u_1=1$ và $u_{n+1}=u_n+2n-1, \forall n \in \mathbb{N}^*$ có $u_{20}=362$} 
  {\True Cho dãy số $(u_n),$  được xác định $\heva{& u_1=1 \\& u_{n+1}=u_n+(-1)^{2n}}$. Số hạng tổng quát $u_n=n$}
  {\True Cho dãy số $(u_n)$ được xác định bởi  $u_1=0$ và $u_{n+1}=n+u_n\,\forall n\ge 1.$ Giá trị của $u_{218}=23653$}
  \loigiai{
  \begin{itemchoice}
   \itemch Mệnh đề \lq\lq  Cho dãy số $(u_n)$ với $u_n = (-1)^n \cdot \dfrac{n}{n+1}$. Khi đó $u_8=\dfrac{8}{9}$\rq\rq\, đúng.
   \itemch Từ $u_{n+1}=u_n+2n-1,  \forall n \in \mathbb{N}^* \Leftrightarrow u_{n+1}-u_n=2n-1, \forall n \in \mathbb{N}^*$.\\
    Do đó
   \begin{eqnarray*}
    &u_2-u_1&=1\\ 
    &u_3-u_2&=3\\ 
     &\vdots\\
     &u_n-u_{n-1}&=2(n-1)-1=2n-3
   \end{eqnarray*}
   Cộng vế theo  vế ta có $u_n-u_1=1+3+\cdots+2n-3=(n-1)^2 \Rightarrow u_n=(n-1)^2+1$.\\
   Vậy  $u_{20}=19^2+1=362$.\\
   Mệnh đề \lq\lq Dãy số $\left(u_n\right)$  biết $u_1=1$ và $u_{n+1}=u_n+2n-1, \forall n \in \mathbb{N}^*$ có $u_{20}=362$\rq\rq\, đúng.
   \itemch Mệnh đề \lq\lq Cho dãy số $(u_n),$ được xác định $\heva{& u_1=1 \\& u_{n+1}=u_n+(-1)^{2n}}$. Số hạng tổng quát $u_n=n$\rq\rq\,  đúng.
   \itemch Từ giả  thiết, ta có $u_{n+1}-u_n=n\,\forall n\ge 1$. \\
    Suy ra $u_{218}=(u_{218}-u_{217})+(u_{217}-u_{216})+\ldots+(u_2-u_1)+u_1=217+216+\cdots+1=\dfrac{217\cdot 218}2=23653$.\\
   Mệnh  đề \lq\lq Cho dãy số $(u_n)$ được xác định bởi $u_1=0$ và $u_{n+1}=n+u_n\,\forall n\ge 1.$ Giá trị của $u_{218}=23653$\rq\rq\, đúng.
  \end{itemchoice}
 } 
\end{ex}

% nguồn: Ok_dungsai_lop11_Lan1
\begin{ex}%[1D2H1-5]%[ma:def8c84f]
 Cho dãy số $(u_n)$  biết $u_n=2^n+1$. Khi đó
 \choiceTF 
  {\True Số hạng thứ $6$ của dãy số  là 65}
  {Số hạng thứ $n+2$ của dãy số  là $u_{n+2}=2\cdot2^n$}
  {\True Dãy số  $(u_n)$ là một dãy số tăng}
  {Dãy  số $(u_n)$ là dãy số bị chặn}
 \loigiai{ 
  \begin{itemchoice}
   \itemch Đúng. Ta có $u_6=2^6+1=65$. 
   \itemch Sai. Ta có  $ u_n=2^n\Rightarrow u_{n+2}=2^{n+2}+1=4\cdot2^n+1$.
   \itemch Đúng. Ta có $u_{n+1}-u_n=2^{n+1}+1-2^n-1=2\cdot2^n-2^n=2^n>0$, $\forall n\in \mathbb{N}^*$. Do đó dãy số $(u_n)$ là một dãy  số tăng.
   \itemch  Sai. Vì dãy số $(u_n)$ không bị chặn trên nên dãy số không bị chặn.
  \end{itemchoice}
  }
\end{ex}

% nguồn: Ok_dungsai_lop11_Lan1
\begin{ex}%[1D2H1-5]%[ma:df7d4905]
 Dãy số có các số hạng cho bởi: $0;\dfrac{1}{2};\dfrac{2}{3};\dfrac{3}{4};\dfrac{4}{5};\ldots $. Các mệnh đề sau đúng hay  sai?
  \choiceTF
  {Số hạng tổng quát của dãy là $u_n=\dfrac{n}{n+1}, \forall  n\in\mathbb{N}$}
  {\True Số hạng tổng quát của dãy là  $u_n=\dfrac{n-1}{n}, \forall n\in\mathbb{N}$}
  {\True Dãy  số tăng}
  {\True Dãy số bị chặn  trên bởi $1$}
  \loigiai{
  \begin{itemchoice}
   \itemch Mệnh đề \lq\lq Số hạng tổng quát của dãy là $u_n=\dfrac{n}{n+1}, \forall  n\in\mathbb{N}$\rq\rq\, sai.
   \itemch Mệnh đề \lq\lq Số hạng tổng quát của dãy là $u_n=\dfrac{n-1}{n}, \forall n\in\mathbb{N}$\rq\rq\,  đúng.
   \itemch Mệnh đề \lq\lq Dãy số tăng\rq\rq\,  đúng.
   \itemch  Mệnh đề \lq\lq Dãy số bị chặn trên bởi $1$\rq\rq\, đúng.
  \end{itemchoice}
 } 
\end{ex}

% nguồn: Ok_dungsai_lop11_Lan1
\begin{ex}%[1D2H1-5]%[ma:fb864715]
 Cho dãy số $\left( u_n\right)$ với $u_n=2n-1 $. Các mệnh đề  sau đúng hay sai?
 \choiceTF 
   {Dãy số bị chặn dưới nhưng không bị chặn trên}
  {\True  Dãy số tăng}
  {Dãy số vừa chặn dưới  vừa bị chặn trên}
  {Dãy  số giảm}
  \loigiai{
  Ta có $u_{n+1}-u_n=\left[2(n+1)-1\right] -\left[2n-1  \right]=(2n+1)-(2n-1)=2>0 $, $\forall n\in \mathbb{N}^*$ nên dãy số đã cho là dãy số tăng, bị chặn dưới bởi $u_1=2\cdot 1-1=1$ và không bị chặn trên.
  \begin{itemchoice}
   \itemch Mệnh đề \lq\lq Dãy số  bị chặn dưới nhưng không bị chặn trên\rq\rq\, sai.
   \itemch Mệnh đề \lq\lq Dãy  số tăng\rq\rq\, đúng.
   \itemch Mệnh đề \lq\lq Dãy số vừa chặn dưới  vừa bị chặn trên\rq\rq\, đúng.
   \itemch Mệnh đề  \lq\lq Dãy số giảm\rq\rq\, sai.
  \end{itemchoice}
  }
\end{ex}

% nguồn: Ok_dungsai_lop11_Lan1
\begin{ex}%[1D2H1-5]%[ma:c5b5e427]
 Cho dãy  số $(u_n)$ biết $u_n=\dfrac{2n-13}{3n-2}$.
  \choiceTF
  {\True Dãy số $(u_n)$  có số hạng thứ mười là $\dfrac{1}{4}$}
  {Dãy số $(u_n)$ là  dãy không tăng, không giảm}
  {\True Dãy số $(u_n)$ là dãy bị chặn} 
  {Dãy số $(u_n)$ bị chặn trên bởi $\dfrac{1}{3}$} 
  \loigiai{
  \begin{itemchoice}
   \itemch $u_{10}=\dfrac{2\cdot 10-13}{3\cdot 10-2}=\dfrac{1}{4}$. 
   \itemch Ta có $u_{n+1}-u_n=\dfrac{2n-11}{3n+1}-\dfrac{2n-13}{3n-2}=\dfrac{34}{(3n+1)(3n-2)}>0$ với  mọi $n\ge 1$.
   \itemch Ta có $u_1=-\dfrac{9}{4}$. Suy ra $u_{n+1}>u_n\forall n\ge 1$ dãy $(u_n)$ là dãy tăng nên dãy bị chặn  dưới bởi $u_1=-\dfrac{9}{4}$.\\
   Mặt  khác $u_n=\dfrac{2}{3}-\dfrac{35}{3(3n-2)}\Rightarrow-\dfrac{9}{4}\le u_n <\dfrac{2}{3}\forall n\ge 1$.\\
   Vậy dãy $(u_n)$  là dãy bị chặn.
   \itemch Dãy bị  chặn trên bởi $-\dfrac{2}{3}$.
  \end{itemchoice}
 } 
\end{ex}

% nguồn: Ok_dungsai_lop11_Lan1
\begin{ex}%[1D2H1-5]%[ma:153dc896]
 Cho  dãy số $(u_n)$, với $u_n=\sin \dfrac{\pi}{n+1}$. Các mệnh đề sau đúng hay sai?
 \choiceTF 
  {Dãy số $(u_n)$ không tăng không  giảm}
  {\True Dãy  số $(u_n)$ là dãy số bị chặn}
  {Dãy số $(u_n)$ là một dãy số tăng} 
  {Số hạng thứ $n+1$ của dãy là  $u_{n+1}=\sin \dfrac{\pi}{n+1}$}
 \loigiai{ 
  Ta có 
   $u_n=\sin \dfrac{\pi}{n+1}\Rightarrow u_{n+1}=\sin \dfrac{\pi }{(n+1)+1}=\sin \dfrac{\pi }{n+2}$.\\
  $u_n=\sin \dfrac{\pi }{n+1}\Rightarrow  -1\le u_n\le 1\Rightarrow (u_n)$ là dãy số bị chặn. \\
  $u_{n+1}-u_n=\sin \dfrac{\pi }{n+2}-\sin \dfrac{\pi }{n+1}<0\,\,\left(0<\dfrac{\pi }{n+2}<\dfrac{\pi }{n+1}\le \dfrac{\pi  }{2}\right)\Rightarrow $ dãy số $(u_n)$ giảm.
  \begin{itemchoice}
   \itemch Mệnh đề \lq\lq Dãy số  $(u_n)$ không tăng không giả\rq\rq\, sai.
   \itemch Mệnh đề \lq\lq Dãy số $(u_n)$ là  dãy số bị chặn\rq\rq\, đúng.
   \itemch Mệnh đề \lq\lq Dãy số $(u_n)$ là một dãy số  tăng\rq\rq\, sai.
   \itemch Mệnh đề \lq\lq ố hạng thứ $n+1$ của dãy là  $u_{n+1}=\sin \dfrac{\pi}{n+1}$\rq\rq\, sai.
  \end{itemchoice}
  }
\end{ex}

% nguồn: Ok_dungsai_lop11_Lan1
\begin{ex}%[1D2H1-5]%[ma:44489703]
 Cho dãy số $\left(u_n\right)$ xác định  bởi $u_n=\sin \dfrac{n\pi}{3}$, với $n \ge 1$. Các mệnh đề sau đúng hay sai?
  \choiceTF
  {\True Số hạng thứ $2$ của dãy  số là $u_2=\dfrac{\sqrt{3}}{2}$}
  {Dãy số $\left(u_n\right)$ là dãy số tăng} 
  {\True Dãy  số $\left(u_n\right)$ bị chặn}
  {Dãy số $\left(u_n\right)$  là dãy số giảm}
  \loigiai{
  Ta có, dãy số $u_n=\sin  \dfrac{n\pi}{3}$ có $u_2=\dfrac{\sqrt{3}}{2}$, dãy không tăng không giảm, $-1\le \sin \dfrac{n\pi }{3}\le 1$. Suy ra dãy số $\left(u_n\right)$ bị chặn.
  \begin{itemchoice}
   \itemch Mệnh đề \lq\lq Số hạng  thứ $2$ của dãy số là $u_2=\dfrac{\sqrt{3}}{2}$\rq\rq\, đúng.
   \itemch Mệnh đề \lq\lq Dãy số $\left(u_n\right)$  là dãy số tăng\rq\rq\, sai.
   \itemch Mệnh đề \lq\lq Dãy  số $\left(u_n\right)$ bị chặn\rq\rq\, đúng.
   \itemch Mệnh đề \lq\lq Dãy số $\left(u_n\right)$ là dãy số  giảm\rq\rq\, sai.
  \end{itemchoice}
  }
\end{ex}

% nguồn: Ok_dungsai_lop11_Lan1
\begin{ex}%[1D2H1-5]%[ma:6835c74d]
 Cho dãy số $(u_n)$, biết $u_n=(-1)^n\cdot 5^{2n+5}$. Các mệnh đề sau đúng hay sai? 
  \choiceTF
  {Dãy số $(u_n)$ bị chặn  dưới và không bị chặn trên}
  {Dãy  số $(u_n)$ bị chặn}
  {Dãy số $(u_n)$ bị  chặn trên và không bị chặn dưới}
  {\True Dãy số $(u_n)$ không  bị chặn}
 \loigiai{ 
  Nếu  $n$ chẵn thì ${u_n}={{5}^{2n+1}}>0$ tăng lên vô hạn (dương vô cùng) khi $n$ tăng lên vô hạn nên dãy $(u_n)$ không bị chặn trên.\\
  Nếu $n$ lẻ thì ${u_n}=-{{5}^{2n+1}}<0$ giảm xuống vô hạn (âm vô cùng) khi $n$ tăng lên vô hạn nên dãy  $(u_n)$ không bị chặn dưới.\\
  Vậy dãy số đã cho không  bị chặn.
  \begin{itemchoice}
   \itemch Mệnh đề \lq\lq Dãy số $(u_n)$  bị chặn dưới và không bị chặn trên\rq\rq\, sai.
   \itemch Mệnh đề \lq\lq Dãy  số $(u_n)$ bị chặn\rq\rq\, sai.
   \itemch Mệnh đề \lq\lq Dãy số $(u_n)$ bị chặn  trên và không bị chặn dưới\rq\rq\, sai.
   \itemch Mệnh đề \lq\lq Dãy số $(u_n)$ không bị chặn\rq\rq\,  đúng.
  \end{itemchoice}
  }
\end{ex}

% nguồn: Ok_dungsai_lop11_Lan1
\begin{ex}%[1D2H1-5]%[ma:2b65d474]
 Cho dãy số $\left(u_n  \right) $với $u_n=\dfrac{1}{n^2+n}$. Các mệnh đề sau đúng hay sai?
  \choiceTF
   {\True $\left( u_n\right) $ bị chặn bởi $\dfrac{1}{2}$}
  {\True $\left(  u_n\right) $ là dãy số giảm}
  {$\left( u_n\right) $ là dãy  số tăng}
  {\True $5$  số hạng đầu của dãy là $\dfrac{1}{2}$, $\dfrac{1}{6}$, $\dfrac{1}{12}$, $\dfrac{1}{20}$, $\dfrac{1}{30}$}
  \loigiai{
  Ta có $u_n>0$  với mọi $n\in \mathbb{N}^*$ và $u_n=\dfrac{1}{n^2+n}=\dfrac{1}{n(n+1)}$ nên
   $$\dfrac{u_{n+1}}{u_n}=\dfrac{\dfrac{1}{(n+1)(n+2)}}{\dfrac{1}{n(n+1)}}=\dfrac{n(n+1)}{(n+1)(n+2)}=\dfrac{n}{n+2}<1.$$
  Do đó, dãy số đã cho  là dãy số giảm và $u_n\le u_1=\dfrac{1}{2}$ với mọi $n\in \mathbb{N}^*$ .\\
  Ta có $5$ số  hạng đầu của dãy là
   \[\dfrac{1}{2},\,\dfrac{1}{6},\,\dfrac{1}{12},\,\dfrac{1}{20},\,\dfrac{1}{30}.\]
  \begin{itemchoice}
   \itemch Mệnh đề \lq\lq $\left( u_n\right) $ bị chặn bởi $\dfrac{1}{2}$\rq\rq\, đúng. 
   \itemch Mệnh đề  \lq\lq $\left( u_n\right) $ là dãy số giảm\rq\rq\, đúng.
   \itemch Mệnh đề \lq\lq $\left( u_n\right)  $ là dãy số tăng\rq\rq\, sai.
   \itemch Mệnh đề \lq\lq $5$  số hạng đầu của dãy là $\dfrac{1}{2}$, $\dfrac{1}{6}$, $\dfrac{1}{12}$, $\dfrac{1}{20}$, $\dfrac{1}{30}$\rq\rq\, đúng.
  \end{itemchoice}
  }
\end{ex}

% nguồn: Ok_dungsai_lop11_Lan2
\begin{ex}%[1D2H1-5]%[ma:61dd248c]
	Cho dãy số $(u_n)$ với $u_n=n^2+n+1$. 
	 \choiceTF
		 {$u_4-u_1=10$}
		{\True $(u_n)$ là dãy  số tăng}
		{$(u_n)$ là dãy số bị chặn trên} 
		{\True  $(u_n)$ là dãy số bị chặn dưới}
	\loigiai{ 
		\begin{itemchoice}
			 \itemch \textbf{Sai.}\\
			Ta  có $u_4-u_1=21-3=18$.
			\itemch  \textbf{Đúng.}\\
			Ta có 
			\begin{align*}
				u_{n+1}-u_n&=\left[\left( n+1 \right)^2+\left( n+1 \right)+1 \right]-\left( n^2+n+1 \right)=\left(  n^2+3n+3 \right)-\left( n^2+n+1 \right)\\
				&=2n+2>0,\ \forall  n\in\mathbb{N}^*
			\end{align*}
			nên $(u_n)$ là dãy số  tăng.
			 \itemch \textbf{Sai.}\\
			Ta có $(u_n)$ không bị  chặn trên.
			\itemch \textbf{Đúng.}\\ 
			Ta có $u_n\ge 3$, $\forall n\in \mathbb{N}^*$ nên $(u_n)$ là dãy số bị  chặn dưới.
		\end{itemchoice}
	} 
\end{ex}

% nguồn: Ok_dungsai_lop11_Lan2
\begin{ex}%[1D2H1-5]%[ma:08bcdbbe]
	Cho  dãy số $(u_n)$ với $u_n=\dfrac{n}{2n+1}$.
	 \choiceTF
		 {\True $u_5=\dfrac{5}{11}$}
		{$\left(u_n\right)$ là dãy số giảm} 
		{$(u_n)$  là dãy số không bị chặn trên}
		{\True $\left(u_n\right)$ là dãy số bị chặn  dưới}
	\loigiai{ 
		\begin{itemchoice}
			 \itemch \textbf{Đúng.}\\
			Ta  có $u_5=\dfrac{5}{2\cdot5+1}=\dfrac{5}{11}$.
			\itemch \textbf{Sai.}\\ 
			 Ta có
			$$\begin{aligned}
			u_{n+1}-u_n&=\dfrac{n+1}{2n+3}-\dfrac{n}{2n+1}=\dfrac{\left( n+1 \right)\left( 2n+1 \right)-n\left( 2n+3 \right)}{\left(  2n+3 \right)\left( 2n+1 \right)}\\
			&=\dfrac{2n^2+3n+1-\left( 2n^2+3n \right)}{\left( 2n+3 \right)\left( 2n+1 \right)}=\dfrac1{\left( 2n+3 \right)\left( 2n+1 \right)}>0, \ \forall  n\in\mathbb{N}^*
		\end{aligned}$$
		nên $(u_n)$ là dãy số  tăng.
		\itemch  \textbf{Sai.}\\
		Ta có  $u_n=\dfrac{n}{2n+1}=\dfrac1{2}-\dfrac1{2\left( 2n+1 \right)}<\dfrac1{2},\ \forall n\in\mathbb{N}^*$ nên $(u_n)$ bị chặn trên.
		\itemch  \textbf{Đúng.}\\
		Ta có $(u_n)$ là  dãy số tăng nên $u_n\ge u_1=\dfrac1{3}$, $\forall n\in\mathbb{N}^*$.
	\end{itemchoice}
 }
\end{ex}

% nguồn: Ok_dungsai_lop11_Lan2
\begin{ex}%[1D2H1-5]%[ma:df805884]
	Cho dãy  số $(u_n)$ với $u_n=\dfrac{an+2}{n+1}$, với $ a$ là số thực.
	 \choiceTF
		{Nếu $u_2=\dfrac{8}{3}$  thì $ a=2$}
		{\True Nếu $ a=1$ thì $(u_n)$ là dãy  số bị chặn}
		{$(u_n)$ là dãy số giảm khi và chỉ khi  $ a<1$}
		{\True $\left(u_n\right)$ là dãy số tăng  khi và chỉ khi $ a>2$}
	\loigiai{ 
		\begin{itemchoice}
			\itemch \textbf{Sai.}\\ 
			Ta có $u_2=\dfrac{8}{3}\Leftrightarrow  \dfrac{2a+2}{3}=\dfrac{8}{3}\Leftrightarrow a=3$.
			\itemch  \textbf{Đúng.}\\
			Ta có với $ a=1$ thì  $u_n=\dfrac{n+2}{n+1}=1+\dfrac1{n+1}$.\\
			Dễ thấy $1<u_n\le 1+\dfrac1{2}=\dfrac{3}{2}$, $\forall n\in \mathbb{N}^*$ nên $(u_n)$ là dãy số  bị chặn.
			\itemch  \textbf{Sai.}\\
			Ta có 
			\begin{align*}
				H&=u_{n+1}-u_n=\dfrac{a\left( n+1 \right)+2}{n+2}-\dfrac{an+2}{n+1}=\dfrac{\left( an+a+2 \right)\left(  n+1 \right)-\left( an+2 \right)\left( n+2 \right)}{\left( n+1 \right)\left( n+2 \right)}\\
				&=\dfrac{a-2}{\left(  n+1 \right)\left( n+2 \right)}.
			\end{align*}
			Ta có $(u_n)$ là dãy số giảm $\Leftrightarrow H<0$, $\forall n\in \mathbb{N}^*\Leftrightarrow  a-2<0\Leftrightarrow a<2$.
			 \itemch \textbf{Đúng.}\\
			Ta có $(u_n)$ là dãy số tăng $\Leftrightarrow H>0$, $\forall n\in \mathbb{N}^*\Leftrightarrow a-2>0\Leftrightarrow  a>2$.
		\end{itemchoice}
	} 
\end{ex}

% nguồn: Ok_dungsai_lop11_Lan1
\begin{ex}%[1D2V1-5]%[ma:075b14db]
 Cho  dãy số $(u_n)$ biết $u_n=\dfrac{2n+1}{n+2}$, $n\in\mathbb{N}^*$.
 \choiceTF 
  {Hai số  hạng đầu lần lượt là $u_1=1$; ${u_2}=\dfrac{5}{3}$}
  {\True Dãy số $(u_n)$ là một dãy số tăng} 
   {Dãy số $(u_n)$ bị chặn trên bởi $\dfrac{33}{17}$}
  {\True Dãy số $(u_n)$ có  duy nhất một số hạng nguyên}
 \loigiai{ 
  \begin{itemchoice}
   \itemch \textbf{Sai}.\\ 
   Ta  có $u_1=\dfrac{2\cdot 1+1}{1+3}=1$; ${u_2}=\dfrac{2\cdot 2+1}{2+2}=\dfrac{5}{4}\ne\dfrac{5}{3}$.
   \itemch \textbf{Đúng}.\\ 
   Ta có $u_n=2-\dfrac{3}{n+2}$ nên khi $n$ tăng thì $\dfrac{3}{n+2}$ giảm  dẫn đến $u_n=2-\dfrac{3}{n+2}$ tăng.\\
   Cho nên  dãy số $(u_n)$ là một dãy số tăng.
    \itemch \textbf{Sai}.\\
   Ta có $u_1\le u_n =2-\dfrac{3}{n+2}<2$ do  $\dfrac{3}{n+2}>0$ nên $u_n=2-\dfrac{3}{n+2}<2$.\\
   Vậy dãy  số $(u_n)$ bị chặn trên bởi $2$.
   \itemch \textbf{Đúng}.\\ 
   Ta có $u_n=2-\dfrac{3}{n+2}$ nên $u_n$ nguyên khi  $3\,\vdots\,(n+2)$ nghĩa là $\hoac{& n+2=\pm 3\\& n+2=\pm 1}\Leftrightarrow \hoac{&n=1\\&n=-5\\&n=-1\\&n=-3.}$\\
   Kết hợp với $n\in\mathbb{N}^*$ ta tìm được $n=1$.\\ 
   Vậy $(u_n)$ chỉ có một số  hạng nguyên là $u_1=1$.
  \end{itemchoice}
  }
\end{ex}

% nguồn: Ok_dungsai_lop11_Lan1
\begin{ex}%[1D2V1-5]%[ma:977a177e]
 Cho dãy số $(u_n)$ biết  $u_n=\dfrac{2n+1}{n+2}, n\in\mathbb{N}^*$. Các mệnh đề sau đây đúng hay sai?
  \choiceTF
  {Hai  số hạng đầu lần lượt là $u_1=1$; ${u_2}=\dfrac{5}{3}$}
   {\True Dãy số $(u_n)$ là một dãy số tăng}
  {Dãy số $(u_n)$ bị chặn  trên bởi $\dfrac{33}{17}$}
   {\True Dãy số $(u_n)$ có duy nhất một số hạng nguyên}
 \loigiai{ 
  \begin{itemchoice}
   \itemch Ta có $u_1=\dfrac{2\cdot 1+1}{1+3}=1$; ${u_2}=\dfrac{2\cdot  2+1}{2+2}=\dfrac{5}{4}\ne\dfrac{5}{3}$.
   \itemch Ta có $u_n=2-\dfrac{3}{n+2}$  nên khi $n$ tăng thì $\dfrac{3}{n+2}$ giảm dẫn đến $u_n=2-\dfrac{3}{n+2}$ tăng.\\
   Cho nên dãy số $(u_n)$  là một dãy số tăng.
   \itemch Ta có $u_1\le u_n  =2-\dfrac{3}{n+2}<2$ do $\dfrac{3}{n+2}>0$ nên $u_n=2-\dfrac{3}{n+2}<2$.\\
   Vậy dãy số $(u_n)$ bị  chặn trên bởi $2$.
   \itemch Ta có $u_n=2-\dfrac{3}{n+2}$ nên $u_n$ nguyên  khi $3\,\vdots\,(n+2)$ nghĩa là $\hoac{& n+2=\pm 3\\& n+2=\pm 1}\Leftrightarrow \hoac{&n=1\\&n=-5\\&n=-1\\&n=-3.}$\\
   Kết hợp với $n\in\mathbb{N}^*$ ta  tìm được $n=1$.\\
   Vậy $(u_n)$ chỉ có một số  hạng nguyên là $u_1=1$.
  \end{itemchoice}
 } 
\end{ex}

% nguồn: Ok_dungsai_lop11_Lan2
\begin{ex}%[1D2V1-5]%[ma:f6435199]
	Cho dãy số $(u_n)$ có số hạng tổng quát  $u_n=\dfrac{n+1}{n+2}$.
	 \choiceTF
		 {\True $u_{n+1}-u_n=\dfrac1{(n+3)(n+2)}$}
		{$u_{n+1}<u_n,\forall  n\in \mathbb{N}^*$}
		{Dãy số $(u_n)$ là dãy số  giảm}
		 {\True Dãy $\left(u_n\right)$ là dãy số bị chặn}
	\loigiai{ 
		\begin{itemchoice}
			 \itemch \textbf{Đúng.}\\
			Xét $u_{n+1}-u_n=\dfrac{n+2}{n+3}-\dfrac{n+1}{n+2}=\dfrac{n^2+4  n+4-\left(n^2+4 n+3\right)}{(n+3)(n+2)}=\dfrac1{(n+3)(n+2)}>0$.
			\itemch \textbf{Sai.}\\ 
			 Ta có $u_{n+1}-u_n=\dfrac1{(n+3)(n+2)}>0$, $\forall n \in \mathbb{N}^*$.\\
			Suy ra $u_{n+1}>u_n$,  $\forall n \in \mathbb{N}^*$.
			\itemch  \textbf{Sai.}\\
			Vì $u_{n+1}>u_n$, $\forall n \in  \mathbb{N}^*$ nên dãy số $(u_n)$ là dãy số tăng.
			\itemch \textbf{Đúng.}\\ 
			Ta có $u_n=\dfrac{n+1}{n+2}>0$, $\forall n  \in \mathbb{N}^*$.\\
			Mặt khác, $u_n=\dfrac{n+1}{n+2}=\dfrac{(n+2)-1}{n+2}=1-\dfrac1{n+2}<1$, $\forall  n \in \mathbb{N}^*$.\\
			Do đó $0<u_n<1, \forall n \in \mathbb{N}^*$ nên dãy $(u_n)$ là dãy số  bị chặn.
		\end{itemchoice}
	} 
\end{ex}

% nguồn: Ok_dungsai_lop11_Lan2
\begin{ex}%[1D2V1-5]%[ma:85d51d8d]
	Cho  dãy số $(u_n)$ có số hạng tổng quát $u_n=n+\dfrac1{n}$.
	 \choiceTF
		{\True  $u_{n+1}>u_n,\forall n\in \mathbb{N}^*$}
		{\True Dãy số $(u_n)$  là dãy số tăng}
		{$u_n\ge 1,\forall  n\in \mathbb{N}^*$}
		{Dãy số đã cho bị chặn  trên}
	\loigiai{ 
		\begin{itemchoice}
			 \itemch \textbf{Đúng.}\\
			Với mọi  số nguyên dương $n$, ta có
			\begin{align*}
				u_{n+1}-u_n &=n+1+\dfrac1{n+1}-\left( n+\dfrac1{n} \right) \\ 
				&  =1-\dfrac1{(n+1)n}=\dfrac{(n+1)n-1}{(n+1)n}>0
			\end{align*}
			(vì $(n+1)n>1$,  $\forall n\ge 1$).
			Suy ra $u_{n+1}>u_n$, $\forall n  \in \mathbb{N}^*$.
			\itemch  \textbf{Đúng.}\\
			Vì  $u_{n+1}>u_n$, $\forall n \in \mathbb{N}^*$ nên dãy số $(u_n)$ là dãy số tăng.
			\itemch  \textbf{Sai.}\\
			 Vì $(u_n)$ là dãy số tăng nên $u_n \ge u_1 =2$, với mọi $n$.\\
			Mặt khác, áp dụng bất  đẳng thức Cô-si cho hai số dương $n$, $\dfrac1{n}$, ta được
			$$n+\dfrac1{n} \geq 2 \sqrt{n  \cdot \dfrac1{n}}=2 \text { hay } u_n \geq 2, \forall n \in \mathbb{N}^* .$$
			\itemch \textbf{Sai.}\\ 
			Dãy số đã cho không bị chặn  trên.
		\end{itemchoice}
	} 
\end{ex}


%%%%%%%%%%%%%%%% TRẢ LỜI NGẮN %%%%%%%%%%%%%%%%

% nguồn: Ok_TraLoiNgan_Lop11_Lan1
\begin{ex}%[1D2V1-5]%[ma:0836199d]
	Cho dãy số $(u_n)$, biết $u_n=n+\dfrac{1}{n}$.  Tìm $m$ để dãy số $(u_n)$ bị chặn dưới bởi $m$.
	 \shortans{2}
	 \loigiai{
	Xét $u_n=n+\dfrac{1}{n}\ge 2\sqrt{n\cdot\dfrac{1}{n}}=2$. (Áp dụng BĐT Côsi)} 
\end{ex}

% nguồn: Ok_TraLoiNgan_Lop11_Lan1
\begin{ex}%[1D2V1-5]%[ma:6ea6a87e]
	Dãy số $\heva{&u_{1}=\sqrt{2}\\&u_{n+1}=\sqrt{u_{n}+2}}$ bị  chặn trên bởi $a$. Tìm $a$.
	 \shortans{2}
	 \loigiai{
		Ta có  $u_{n}>1$.\\
		Giả sử tồn tại  $u_{n} \geq 2 \Rightarrow \sqrt{u_{n+1}+2} \geq 2 \Rightarrow u_{n-1} \geq 2$.\\
		Nếu tồn tại $u_{n} \geq 2$ thì suy ra $u_{n-1} \geq 2$, từ đó suy ra được $u_{n-2}, u_{n-3}, \ldots, u_{2}, u_{1} \geq 2$ vô lý.\\ 
		Do $u_{1}<2$. Nên điều giả sử là sai. Suy  ra $u_{n}<2$.\\
		 Xét $u_{n+1}-u_{n}=\sqrt{u_{n}+2}-u_{n}=\dfrac{u_{n}+2-u_{n}^2}{\sqrt{u_{n}+2}+u_{n}}=\dfrac{(2-u_{n})(1+u_{n})}{\sqrt{u_{n}+2}+u_{n}}>0$.\\
		Suy ra: $u_{n+1}>u_{n}$ nên $(u_{n})$ là  dãy tăng.\\
	Vậy dãy đã  cho tăng và bị chặn trên bởi $2$.}
\end{ex}

% nguồn: Ok_TraLoiNgan_Lop11_Lan2
\begin{ex}%[1D2V1-5]%[ma:53f76ac3]
	Dãy số $\heva{&u_1=\sqrt{2}\\&u_{n+1}=\sqrt{u_n+2}}$ bị chặn trên bởi $a$. Khi đó $a$ có giá trị bằng  bao nhiêu?
	\shortans{2} 
	 \loigiai{
		Với $n=1$, ta có $u_1  = \sqrt{2} < 2$ (hiển nhiên)\\
		Giả sử $n = k$ thì $u_k  < 2 $ là đúng.\\
		Ta chứng minh: với  $n = k + 1$ thì $u_{k+1} < 2$.\\
		 Thật vậy ta có\\
		 $u_{k+1} = \sqrt{u_k+2}$\\
		Mà $u_k < 2 $ nên $\sqrt{u_k + 2} <  \sqrt{2+2} = 2$\\
		Suy ra $u_{k+1} <  2$.\\
		Vậy $u_n < 2$ với mọi $n \in \mathbb{N}^*$. 
	 }
\end{ex}


%%%%%%%%%%%%%%%% TỰ LUẬN %%%%%%%%%%%%%%%%

% nguồn: Ok_TuanLuan_Toan11_Lan2
\begin{ex}%[1D2H1-5]%[ma:0b5dc59e]
	Xét tính tăng, giảm và bị chặn  của dãy số $(u_n)$ với
	\begin{listEX}[2]
		\item $u_n = \dfrac{n+1}{n+2}$. 
		 \item $u_n = \dfrac{1-n}{n+1}$.
	\end{listEX}
	 \loigiai{
		\begin{listEX}
			\item Ta có \begin{eqnarray*}
			u_{n+1}-u_n &=& \dfrac{n+2}{n+3}  - \dfrac{n+1}{n+2}\\
			&=& \dfrac{(n+2)^2-(n+3)(n+1)}{(n+2)(n+3)}\\ 
			&=& \dfrac{1}{(n+2)(n+3)} >  0, \forall n \in \mathbb{N}^*.
		\end{eqnarray*}
		Suy ra $u_{n+1} > u_n, \forall n \in  \mathbb{N}^*$.\\
		Vậy $(u_n)$ là dãy  số tăng.\\
		 Mặt khác ta có $u_n = \dfrac{n+1}{n+2} > 0$, $\forall n \in \mathbb{N}^*$ suy ra $(u_n)$ bị chặn dưới.\\
		$u_n=1-\dfrac{1}{n+2}<1$, $\forall n \in  \mathbb{N}^*$, suy ra $(u_n)$ bị chặn trên.\\
		Ta thấy dãy số $(u_n)$ vừa bị chặn trên vừa bị chặn dưới, suy ra  dãy số $(u_n)$ bị chặn.\\
		Vậy $(u_n)$ là dãy số  tăng và bị chặn.
		\item Ta có \begin{eqnarray*}
		u_{n+1}-u_n &=& \dfrac{1-(n+1)}{n+2} -  \dfrac{1-n}{n+1}\\
		&=& \dfrac{-n}{n+2} -  \dfrac{1-n}{n+1}\\
		&=& \dfrac{-n(n+1)-(1-n)(n+2)}{(n+2)(n+1)} \\ 
		&=& \dfrac{2}{(n+2)(n+1)} < 0, \forall n  \ge 1.
	\end{eqnarray*}
	Do  đó $(u_n)$ là dãy số giảm.\\
	Do $u_n = \dfrac{1-n}{n+1} = \dfrac{-(n+1)+2}{n+1} = -1 +  \dfrac{2}{n+1} > -1$, $\forall n \ge 1$, nên dãy $(u_n)$ bị chặn dưới.\\
	Dãy  $(u_n)$ cũng bị chặn trên vì $u_n = -1+\dfrac{2}{n+1} \le -1+\dfrac{2}{1+1} = 0$, $\forall n \ge 1$.\\
	Do đó $(u_n)$ là dãy số bị chặn. 
\end{listEX}}
\end{ex}

% nguồn: Ok_TuanLuan_Toan11_Lan2
\begin{ex}%[1D2H1-5]%[ma:eacb3cad]
	Trong các dãy số $(u_n)$ sau, dãy số nào bị chặn dưới, bị chặn  trên, bị chặn?
	\begin{listEX}[2]
		\item $u_n = n-1$. 
		\item $u_n  = \dfrac{n+1}{n+2}$.
		 \item $u_n = \sin n$.
		\item  $u_n = (-1)^{n-1} n^2$.
	\end{listEX}
	\loigiai{ 
		\begin{listEX}
			\item  Ta có $u_n = n-1 \ge 0$, $(\forall n \in \mathbb{N^*})$ suy ra $u_n$ bị chặn dưới.
			\item Ta có $u_n = \dfrac{n+1}{n+2} = 1 - \dfrac{1}{n+2} < 1$; $u_n = \dfrac{n+1}{n+2} = 1 - \dfrac{1}{n+2} \ge 0$, $(\forall  n \in \mathbb{N}^*)$ suy ra $u_n$ bị chặn.
			\item $u_n = \sin n$  do đó $-1 \le u_n \le 1$, $(\forall n \in \mathbb{N}^*)$ suy ra $u_n$ bị chặn.
			\item  Ta có
			\begin{itemize}
				\item $\lim\limits u_n =\lim\limits (-1)^{n-1}\cdot n^2 =-\infty$  nếu $n$ là số tự nhiên lẻ.
				\item $\lim\limits u_n =\lim\limits (-1)^{n-1}\cdot  n^2 =+\infty$ nếu $n$ là số tự nhiên chẵn.
			\end{itemize}
			Suy ra dãy không bị  chặn trên cũng không bị chặn dưới.
		\end{listEX}
	} 
\end{ex}

% nguồn: Ok_TuanLuan_Toan11_Lan2
\begin{ex}%[1D2H1-5]%[ma:068bdd6e]
	Chứng minh rằng dãy  số $(u_n)$ với $u_n = \dfrac{1}{n^2+n}$ bị chặn.
	\loigiai{ 
		 Ta có $u_n > 0$ với mọi $n \in \mathbb{N}^*$ nên dãy số $(u_n)$ bị chặn dưới.\\
		Do $n^2+n \ge 2$ với  mọi $n \in \mathbb{N}^*$ nên $u_n \le \dfrac{1}{2}$ với mọi $n \in \mathbb{N}^*$ suy ra dãy số $(u_n)$ bị chặn trên.\\
		 Vậy
	dãy số $(u_n)$ bị  chặn.}
\end{ex}

% nguồn: Ok_TuanLuan_Toan11_Lan2
\begin{ex}%[1D2H1-5]%[ma:a8d082b9]
	Xét tính  bị chặn của dãy số $(u_n)$ biết $u_n=\dfrac{4n+5}{n+1}$.
	\loigiai{ 
		Ta có  $u_n=\dfrac{4n+5}{n+1} > 0, \forall n \in \mathbb{N}^*$.\\
		Và $u_n=\dfrac{4n+5}{n+1}=\dfrac{4(n+1)+1}{n+1}=4+\dfrac{1}{n+1} \leq 4+\dfrac{1}{2}=\dfrac{9}{2} \Rightarrow u_n \leq \dfrac{9}{2}, \forall  n \in \mathbb{N}^*.$\\
		 Suy ra $0< u_n \leq \dfrac{9}{2}, \forall n \in \mathbb{N}^*$.\\
		Vậy dãy số $(u_n)$ bị chặn.\\ 
		\textbf{Giải nhanh:} dãy số  $(u_n)$ có $u_n$ có bậc của tử bằng bậc của mẫu nên bị chặn
	 }
\end{ex}

% nguồn: Ok_TuanLuan_Toan11_Lan2
\begin{ex}%[1D2H1-5]%[ma:2dd80024]
	Xét tính bị chặn  của dãy số $(u_n)$ biết $u_n=\dfrac{-1}{2n+3}$.
	\loigiai{ 
		Ta có $2n+3\geq 5, \forall n \in \mathbb{N}^* \Rightarrow 0< \dfrac{1}{2n+3} \leq \dfrac{1}{5}, \forall n \in \mathbb{N}^* \Rightarrow-\dfrac{1}{5} \leq \dfrac{-1}{2n+3} <  0, \forall n \in \mathbb{N}^*$.\\
		 Suy ra $-\dfrac{1}{5} \leq u_n < 0, \forall n \in \mathbb{N}^*$.\\
		Hay  dãy số $(u_n)$ bị chặn.\\
		\textbf{Giải nhanh:} dãy số $(u_n)$ có $u_n$ có bậc của tử thấp  hơn bậc của mẫu nên bị chặn.
	} 
\end{ex}

% nguồn: Ok_TuanLuan_Toan11_Lan2
\begin{ex}%[1D2H1-5]%[ma:124ffb48]
	Xét tính bị chặn của các dãy số $(u_n)$ cho bởi 
	\begin{listEX}[2]
		\item $u_n = \dfrac{1}{1\cdot3} + \dfrac{1}{3\cdot5} +  \cdots + \dfrac{1}{(2n-1)(2n+1)}$.
		 \item $u_n = \dfrac{1}{\sqrt{n^2+1}} + \dfrac{1}{\sqrt{n^2+2}} + \cdots + \dfrac{1}{\sqrt{n^2+n}}$.
	\end{listEX}
	\loigiai{\begin{listEX}
		\item Ta có \begin{eqnarray*}
		u_n  &=& \dfrac{1}{1\cdot3} + \dfrac{1}{3\cdot5} + \cdots + \dfrac{1}{(2n-1)(2n+1)}\\
		&=&  \dfrac{1}{2} \left( \dfrac{1}{1} - \dfrac{1}{3} \right) + \dfrac{1}{2} \left( \dfrac{1}{3} - \dfrac{1}{5} \right) + \cdots + \dfrac{1}{2} \left( \dfrac{1}{2n-1} - \dfrac{1}{2n+1} \right)\\
		 &=& \dfrac{1}{2} \left( \dfrac{1}{1} - \dfrac{1}{2n+1} \right)\\
		&=& \dfrac{1}{2} - \dfrac{1}{2(2n+1)}. 
	\end{eqnarray*}
	Do đó $0 < u_n < \dfrac{1}{2}$,  $\forall n \in \mathbb{N}^*$. \\
	 Vậy $(u_n)$ là dãy số bị chặn.
	\item Ta có $$0 < \dfrac{1}{\sqrt{n^2+1}} + \dfrac{1}{\sqrt{n^2+2}} + \cdots  + \dfrac{1}{\sqrt{n^2+n}} < \dfrac{1}{\sqrt{n^2}} + \dfrac{1}{\sqrt{n^2}} + \cdots + \dfrac{1}{\sqrt{n^2}} = \dfrac{n}{\sqrt{n^2}} = \dfrac{n}{n} = 1, \forall n \in \mathbb{N}^*.$$
	Do  đó $0 < u_n < 1$, $\forall n \in \mathbb{N}^*$.\\
	Vậy $(u_n)$ là dãy số  bị chặn.
\end{listEX}}
\end{ex}

% nguồn: Ok_TuanLuan_Toan11_Lan2
\begin{ex}%[1D2V1-5]%[ma:80f4b4f2]
	 Xét tính bị chặn của dãy số $(u_n)$ biết $u_n=\dfrac{1}{1^2}+\dfrac{1}{2^2}+\dfrac{1}{3^2}+\cdots+\dfrac{1}{n^2}$.
	 \loigiai{
		Xét  $\dfrac{1}{k^2} < \dfrac{1}{(k-1)k}=\dfrac{1}{k-1}-\dfrac{1}{k}$, $\forall k\ge 2$.
		Suy ra $u_n < \dfrac{1}{2}+1-\dfrac{1}{2}+\dfrac{1}{2}-\dfrac{1}{3}+\dfrac{1}{3}-\dfrac{1}{4}+\dfrac{1}{4}-\dfrac{1}{5}+\dfrac{1}{5}-\dfrac{1}{6}+\cdots+\dfrac{1}{n-1}-\dfrac{1}{n}=\dfrac{3}{2}-\dfrac{1}{n} < \dfrac{3}{2}$. 
		 $\Rightarrow u_n\le\dfrac{3}{2}$ và $u_n >0$, $\forall n\in\mathbb{N}^*$
		Vậy $(u_n)$ bị chặn. 
	 }
\end{ex}

% nguồn: Ok_TuanLuan_Toan11_Lan2
\begin{ex}%[1D2V1-5]%[ma:e6beabbd]
	Xét tính bị chặn  của dãy số $(u_n)$ biết $u_n=\dfrac{n^3}{n^2+1}$.
	 \loigiai{
		Ta  có $u_n=\dfrac{n^3}{n^2+1} > 0$, $\forall n\in\mathbb{N}^*\Rightarrow(u_n)$ bị chặn dưới.\\
		Ta có $u_n=\dfrac{n^3}{n^2+1} =n - \dfrac{n}{n^2+1}< n$, $\forall n\in\mathbb{N}^*$ hay không  tồn tại số $M$ nào để $u_n \leq M$ $\forall n\in\mathbb{N}^*$ không bị chặn trên.\\
	} 
\end{ex}
